% concatenated from arxiv_main.tex, macros.tex
\documentclass[11pt]{article}

\usepackage{amsmath,amssymb,amsfonts,amsthm}
\usepackage{thmtools}
\usepackage[numbers,sort]{natbib}
\usepackage{xcolor}
\definecolor{dkred}  {rgb}{.7,0,0} 
\definecolor{midnight}  {rgb}{0,0,.5} 
\definecolor{forest}  {rgb}{0,.6,0} 
\usepackage[colorlinks,
            citecolor=forest,
			linkcolor=midnight,
			anchorcolor=red,
			urlcolor=midnight]{hyperref}  % must be loaded before xurl
\usepackage{xurl}      % enhances \url to break long links cleanly
\newcommand{\becca}[1]{\textcolor{dkred}{[becca: #1]}}
\newcommand{\btodo}[1]{\textcolor{red}{[to do: #1]}}
\newcommand{\ram}[1]{\textcolor{violet}{#1}}

% if you use cleveref..
\usepackage[capitalize,noabbrev]{cleveref}
\Crefname{equation}{Eq.}{Eqs.}

% Recommended, but optional, packages for figures and better typesetting:
\usepackage{microtype}
\usepackage{graphicx}
\usepackage{subcaption}
\usepackage{booktabs} % for professional tables

\usepackage{appendix}
\usepackage{verbatim}




% Attempt to make hyperref and algorithmic work together better:
\newcommand{\theHalgorithm}{\arabic{algorithm}}



\usepackage{arxiv}


\usepackage{amsmath}
\usepackage{amssymb}
\usepackage{mathtools}
\usepackage{amsthm}

% \usepackage{natbib}
\usepackage{xcolor}
\usepackage{graphicx} % Required for inserting images
\usepackage{amsfonts}
\usepackage{amssymb}
\usepackage{amsmath}
\usepackage{amsthm}

\usepackage{subcaption}
\usepackage{bbm, dsfont}


\newcommand{\fitparenth}[1]{\left(#1\right)}
\newcommand{\fitbracket}[1]{\left[#1\right]}
\newcommand{\inprod}[1]{\left\langle#1\right\rangle}

\newcommand{\vect}[1]{\mathbf{#1}}

\newcommand{\RR}{\mathbb{R}}
\newcommand{\NN}{\mathbb{N}}
\newcommand{\CalN}{\mathcal{N}}
\newcommand{\CalS}{\mathcal{S}}
% \newcommand{\expect}[2][]{\mathbb{E}_{#1}\left[#2\right]}
\NewDocumentCommand{\expect}{O{} m}{\mathbb{E}_{#1}\left[#2\right]}

\newcommand{\andrew}[1]{\textcolor{olive}{[andrew: #1]}}
%%%%%%%%%%%%%%%%%%%%%%%%%%%%%%%%
% THEOREMS
%%%%%%%%%%%%%%%%%%%%%%%%%%%%%%%%
\theoremstyle{plain}
\newtheorem{theorem}{Theorem}[section]
\newtheorem{proposition}[theorem]{Proposition}
\newtheorem{lemma}[theorem]{Lemma}
\newtheorem{corollary}[theorem]{Corollary}
% \theoremstyle{definition}
\newtheorem{definition}[theorem]{Definition}
\newtheorem{assumption}[theorem]{Assumption}
% \theoremstyle{remark}
\newtheorem{remark}{Remark}
% \theoremstyle{definition}

\Crefname{proposition}{Proposition}{Propositions}

% ---------- Author/affiliation formatting ----------
\usepackage{authblk}
\setlength{\affilsep}{0.3em}      % space between author block and affiliations
\renewcommand\Authfont{\large}    % author name size
\renewcommand\Affilfont{\small} % affiliation size

\usepackage[compact]{titlesec}

% ---------- Title & authors ----------

\author[*,1,2,3]{Andrew Dennehy}
\author[3]{Ramchandran Muthukumar}
\author[1,2,3,4,5,6]{\\Rebecca Willett}
\author[1,2,5,6]{Nisha Chandramoorthy}

\affil[1]{Committee on Computational and Applied Mathematics, University of Chicago}
\affil[2]{Department of Statistics, University of Chicago}
\affil[3]{Data Science Institute, University of Chicago}
\affil[4]{Department of Computer Science, University of Chicago}
\affil[5]{NSF-Simons National Institute for Theory and Mathematics in Biology}
\affil[6]{NSF-Simons National AI Institute for the Sky}
\affil[*]{Corresponding Author: \href{mailto:nishac@uchicago.edu}{nishac@uchicago.edu}}

% \author{%
%   Andrew Dennehy  \\
%   University of Chicago Committee on Computational and Applied Mathematics, Chicago, IL 60615 \\
%   \texttt{adennehy@uchicago.edu} 
%   \And 
%   Ramchandran Muthukumar  \\
%   University of Chicago Data Science Institute, Chicago, IL 60615 \\
%   \And 
%   Rebecca Willett  \\
%   Department of Statistics and Department of Computer Science \\
%   Committee on Computational and Applied Mathematics\\
%   University of Chicago, Chicago, IL 60615 \\
%   NSF-Simons National Institute for Theory and Mathematics in Biology, Chicago, IL 60611 
%   \And 
%   Nisha Chandramoorthy  \\
%   University of Chicago Committee on Computational and Applied Mathematics, Chicago, IL 60615 \\
%   University of Chicago Department of Statistics, Chicago, IL 60615 \\
%   %\thanks{Use footnote for providing further information
%     % about author (webpage, alternative address)---\emph{not} for acknowledging
%     % funding agencies.}
%   % examples of more authors
%   % \And
%   % Coauthor \\
%   % Affiliation \\
%   % Address \\
%   % \texttt{email} \\
%   % \AND
%   % Coauthor \\
%   % Affiliation \\
%   % Address \\
%   % \texttt{email} \\
%   % \And
%   % Coauthor \\
%   % Affiliation \\
%   % Address \\
%   % \texttt{email} \\
%   % \And
%   % Coauthor \\
%   % Affiliation \\
%   % Address \\
%   % \texttt{email} \\
% }

% Todonotes is useful during development; simply uncomment the next line
%    and comment out the line below the next line to turn off comments
\usepackage[textsize=tiny]{todonotes}
%\usepackage[disable,textsize=tiny]{todonotes}

% The \icmltitle you define below is probably too long as a header.
% Therefore, a short form for the running title is supplied here:
% \icmltitlerunning{Learning Parameters with Diffusion Models}

\ProvidesPackage{latex-macros}

% Require packages
\RequirePackage{amsmath,amssymb,amsthm}
\RequirePackage{bm,bbm}
\RequirePackage{mathtools}
\RequirePackage{hyperref}

% mathbb letters
\newcommand{\Abb}{\mathbb{A}}
\newcommand{\Cbb}{\mathbb{C}}
\newcommand{\Dbb}{\mathbb{D}}
\newcommand{\Ebb}{\mathbb{E}}
\newcommand{\Fbb}{\mathbb{F}}
\newcommand{\Gbb}{\mathbb{G}}
\newcommand{\Hbb}{\mathbb{H}}
\newcommand{\Ibb}{\mathbb{I}}
\newcommand{\Jbb}{\mathbb{J}}
\newcommand{\Kbb}{\mathbb{K}}
\newcommand{\Lbb}{\mathbb{L}}
\newcommand{\Mbb}{\mathbb{M}}
\newcommand{\Nbb}{\mathbb{N}}
\newcommand{\Obb}{\mathbb{O}}
\newcommand{\Pbb}{\mathbb{P}}
\newcommand{\Qbb}{\mathbb{Q}}
\newcommand{\Rbb}{\mathbb{R}}
\newcommand{\Sbb}{\mathbb{S}}
\newcommand{\Tbb}{\mathbb{T}}
\newcommand{\Ubb}{\mathbb{U}}
\newcommand{\Vbb}{\mathbb{V}}
\newcommand{\Wbb}{\mathbb{W}}
\newcommand{\Xbb}{\mathbb{X}}
\newcommand{\Ybb}{\mathbb{Y}}
\newcommand{\Zbb}{\mathbb{Z}}

% mathcal letters
\newcommand{\cA}{\mathcal{A}}
\newcommand{\cB}{\mathcal{B}}
\newcommand{\cC}{\mathcal{C}}
\newcommand{\cD}{\mathcal{D}}
\newcommand{\cE}{\mathcal{E}}
\newcommand{\cF}{\mathcal{F}}
\newcommand{\cG}{\mathcal{G}}
\newcommand{\cH}{\mathcal{H}}
\newcommand{\cI}{\mathcal{I}}
\newcommand{\cJ}{\mathcal{J}}
\newcommand{\cK}{\mathcal{K}}
\newcommand{\cL}{\mathcal{L}}
\newcommand{\cM}{\mathcal{M}}
\newcommand{\cN}{\mathcal{N}}
\newcommand{\cO}{\mathcal{O}}
\newcommand{\cP}{\mathcal{P}}
\newcommand{\cQ}{\mathcal{Q}}
\newcommand{\cR}{\mathcal{R}}
\newcommand{\cS}{\mathcal{S}}
\newcommand{\cT}{\mathcal{T}}
\newcommand{\cU}{\mathcal{U}}
\newcommand{\cV}{\mathcal{V}}
\newcommand{\cW}{\mathcal{W}}
\newcommand{\cX}{\mathcal{X}}
\newcommand{\cY}{\mathcal{Y}}
\newcommand{\cZ}{\mathcal{Z}}

%% var letters
% lower
\newcommand{\veps}{\varepsilon}
\newcommand{\vkap}{\varkappa}
\newcommand{\vphi}{\varphi}
\newcommand{\vpi}{\varpi}
\newcommand{\vrho}{\varrho}
\newcommand{\vthe}{\vartheta}
% upper
\newcommand{\vGam}{\varGamma}
\newcommand{\vDel}{\varDelta}
\newcommand{\vLam}{\varLambda}
\newcommand{\vPsi}{\varPsi}
\newcommand{\vPi}{\varPi}
\newcommand{\vPhi}{\varPhi}
\newcommand{\vSig}{\varSigma}

%% set ops
% delimited
\DeclareMathOperator*{\setI}{\bigcap}
\DeclareMathOperator*{\setU}{\bigcup}
% non-delimited
\newcommand{\setD}{\setminus}
\newcommand{\setDsym}{\bigtriangleup}

% argmin/argmax/arginf/argsup as operators
% with up/down limits
\DeclareMathOperator*{\argmin}{argmin}
\DeclareMathOperator*{\argmax}{argmax}
\DeclareMathOperator*{\arginf}{arginf}
\DeclareMathOperator*{\argsup}{argsup}

% various shorthand macros
\newcommand{\dist}{\mathrm{dist}}
\newcommand{\proj}{\mathrm{proj}}
\newcommand{\prox}{\mathrm{prox}}
\newcommand{\grad}{\nabla}
\newcommand{\conv}{\circledast}
\newcommand{\convhull}{\mathrm{conv}}
\newcommand{\iid}{\overset{\mathrm{i.i.d.}}{\sim}}
% note: span is a latex primitive
\newcommand{\vspan}{\mathrm{span}}
\newcommand{\rank}{\mathrm{rank}}
\newcommand{\tr}{\mathrm{Tr}}
\newcommand{\sign}{\mathrm{sign}}
\newcommand{\range}{\mathrm{range}}
\newcommand{\tdef}{\triangleq}
% middle bar, for use in delimited expressions like sets
\newcommand{\mmid}{\ \middle| \ }
% Positive semidefinite, positive definite cones
\newcommand{\psd}[1]{\Sbb^{#1}_+}
\newcommand{\pd}[1]{\Sbb^{#1}_{++}}
% up/down arrows for limits
\newcommand{\dto}{\downarrow}
\newcommand{\uto}{\uparrow}
\newcommand{\T}{\mathsf{T}}
% unit ball
\newcommand{\ball}{\mathbb{B}}
% multiset
\newcommand{\multo}{\rightrightarrows}
% KL divergence
\newcommand{\dkl}[2]{D_{\mathrm{KL}}\left({#1}\,\|\,{#2}\right)}

% convergence symbols
\newcommand{\asto}{\overset{\mathrm{a.s.}}{\longrightarrow}}
\newcommand{\probto}{\overset{\Pbb}{\longrightarrow}}
\newcommand{\lpto}[1]{\overset{L^{#1}}{\longrightarrow}}
\newcommand{\weakto}{\overset{\mathrm{weak}}{\longrightarrow}}
\newcommand{\normto}{\overset{\mathrm{norm}}{\longrightarrow}}

% convex analysis
\newcommand{\diag}{\mathrm{diag}}
\newcommand{\gph}{\mathrm{gph}}
\newcommand{\epi}{\mathrm{epi}}
\newcommand{\dom}{\mathrm{dom}}
\newcommand{\intr}[1]{\mathrm{int}\left( #1 \right)}
\newcommand{\clos}[1]{\mathrm{cl}\left( #1 \right)}
\newcommand{\relint}[1]{{#1}^{\circ}}

% shorthands for parentheses / brackets
\newcommand{\pars}[1]{\left( #1 \right)}
\newcommand{\brks}[1]{\left[ #1 \right]}

% macros for delimited expressions
\newcommand{\expfun}[1]{\exp\left\{ {#1} \right\}}
\newcommand{\trace}[1]{\tr \left( #1 \right)}
% indicator function
\newcommand{\1}{\mathbbm{1}}

%%%%%%
% The macros below used to be explicit about size, but often required
% manual adjustments. With mathtools, these can be adjusted on the fly:
% Example: $\norm{x}$ vs. $\norm[\Big]{x}$ vs. $\norm*{x}$.
%          (unsized)      (sized as \Big)      (automatically sized)
%%%%%%

% Inner products
\DeclarePairedDelimiter{\ip}{\langle}{\rangle}

% Absolute values
\DeclarePairedDelimiter{\abs}{\lvert}{\rvert}

% Norms
\DeclarePairedDelimiter{\norm}{\lVert}{\rVert}
\DeclarePairedDelimiter{\frobnorm}{\lVert}{\rVert_{\mathsf{F}}}
\DeclarePairedDelimiter{\opnorm}{\lVert}{\rVert_{\mathsf{op}}}

% Set brackets
\DeclarePairedDelimiter{\set}{\lbrace}{\rbrace}

% Floors and ceilings
\DeclarePairedDelimiter{\floor}{\lfloor}{\rfloor}
\DeclarePairedDelimiter{\ceil}{\lceil}{\rceil}

% differential in integrals
\newcommand{\dd}[1]{\mathrm{d}{#1}}

% expectation, probability and variance, with optional subscripts
% \newcommand{\expec}[2][]{
% 	\ifx\\#1\\
% 		\Ebb \left[ #2 \right]
% 	\else
% 		\Ebb_{#1} \left[ #2 \right]
% 	\fi
% }
\newcommand{\prob}[2][]{
	\ifx\\#1\\
		\Pbb \left( #2 \right)
	\else
		\Pbb_{#1} \left( #2 \right)
	\fi
}
\newcommand{\var}[2][]{
	\ifx\\#1\\
		\mathrm{Var}_{#1}\left( #2 \right)
	\else
		\mathrm{Var}\left( #2 \right)
	\fi
}

% partial fractions
\newcommand{\pfrac}[2]{\frac{\partial #1}{\partial #2}}

% matrices
\newcommand{\pmx}[1]{\begin{pmatrix} #1 \end{pmatrix}}
\newcommand{\bmx}[1]{\begin{bmatrix} #1 \end{bmatrix}}
\newcommand{\Bmx}[1]{\begin{Bmatrix} #1 \end{Bmatrix}}


\newcommand{\btheta}{{\boldsymbol{\theta}}}
\newcommand{\vx}{\vect{x}}
\newcommand{\vs}{\vect{s}}
\newcommand{\vy}{\vect{y}}
\newcommand{\vz}{\vect{z}}


\newcommand{\ldsm}{\ell_\mathrm{DSM}}
\newcommand{\lsm}{\ell_\mathrm{SM}}
\newcommand{\lnn}{\ell_{\mathrm{NN}}}

\newcommand{\pdist}[2][]{p_{#1}^{(#2)}}
\newcommand{\qdist}[2][]{q_{#1}^{(#2)}}
\newcommand{\pbare}[1]{p_{#1}} 
\newcommand{\scorebare}[1]{\vect{s}_{#1}} 
\newcommand{\score}[2][]{\vect{s}_{#1}^{(#2)}}
\newcommand{\separation}{\Delta}




\usepackage{titling}
\settowidth{\thanksmarkwidth}{*}
\setlength{\thanksmargin}{-\thanksmarkwidth}
\title{Diffusion models recover accurate mixture weights despite score function insensitivity}
\begin{document}



\maketitle

\begin{abstract}
Score-based generative models exhibit a puzzling behavior: they often appear to cover all modes of a target multimodal distribution and yet may fail to learn the correct relative mode amplitudes, which can be interpreted as mixture weights. 
% Early evidence indicated strong mode coverage, however, recent studies suggest a more nuanced picture with substantial variation based on the training recipe, %neural architecture 
% and the inference-time sampling algorithm. We explain this variation through a rigorous theoretical perspective 
We resolve this apparent paradox 
by relating the diffusion score matching (DSM) loss to the error in estimating mixture weights from generated samples. We show that, even when the target score is insensitive to mixture weights, generated samples can recover the weights accurately if the scores at intermediate noise levels are informative about the weights. Accordingly, we define the \textit{diffusion score sensitivity index} (DSSI) as the variation in the DSM loss relative to changes in a parameter. We then show that the DSSI governs the accuracy with which the parameter of the target distribution can be estimated from generated samples.
For Gaussian mixtures in arbitrary dimensions, we prove that the mixture weight estimation errors are on the same order as the DSM loss under mild conditions. Empirically, we show the emergence of sensitivity during the noising process of benchmark data distributions under typical noise schedules, and that these sensitivity values predict how well a well-trained model recovers mixture weights. Furthermore, we show that the choice of noise schedule can reduce diffusion sensitivity, leading to mode amplification. Although we focus on mixture weights, the proposed sensitivity framework governs the recovery of any qualitative parameter of the target distribution. 
 
\end{abstract}



\section{Introduction}
% \todo[inline]{Affiliations}
Generative models are powerful tools that are increasingly used in scientific research \citep{lanusse2021deep,price2023gencast,han2025learning}, medicine \citep{loni2025review,fahrner2025generative,teo2025generative}, the justice system \citep{ludwig2024machine,posner2025judge}, and more. In these and other settings, users rely on generative models to accurately reflect the qualitative properties of the true target distribution. Past work notes that learning the score of a target distribution may be insufficient for accurate sampling \citep{koehler2022statisticalefficiencyscorematching, chen2021dimensionfreelogsobolevinequalitiesmixture, li2024softmixturedenoisingexpressive}, particularly for canonical multimodal distributions such as Gaussian mixtures \citep{1100034, yakowitz1968identifiability}. 
Generative models that incorrectly reflect the amplitudes or variances of modes \citep{zhang2025collapse, roos2026met, hakemi2025deeper} may hinder uncertainty quantification and the recovery of qualitative parameters \citep{bouchet2019rare, addison2024cpmgem, mudur2023cosmological}. %such as mixture weights of a multimodal data distribution. 
Our work proposes an explicit evaluation of generative models through the lens of parameter recovery. 
%Thus, recovery of qualitative parameters from generated samples is therefore an explicit test of whether learned generative models capture the critical structural properties of the target data distribution.
Such an evaluation augments sample-quality metrics, such as the Frechet inception distance (FID) \citep{NIPS2017_8a1d6947}, which may not explicitly test for target-specific structure, such as multimodality \citep{sajjadi2018assessing}. %,  This and as such complementing 

%In general, if we seek to learn a distribution with key parameters that affect qualitative properties of the distribution (e.g., mixture weights), then those parameters should be accurately reflected in the distribution corresponding to the learned generative model.  
% This insensitivity means that one could potentially learn a model with highly inaccurate relative mode amplitudes and still have a near-optimal value of the loss. 
\textbf{In this work, we present an information-theoretic mechanism that determines when such key parameters may be represented accurately by diffusion models.} In the Gaussian mixture setting, even when the score matching loss is insensitive to mixture weights, learning a diffusion model anneals the score matching loss across steps of the diffusion process \citep{song2019generative}; this annealing leads to sensitivity to mixture weights, resulting in accurate estimates that were unanticipated by past analyses that focused only on the (unannealed) score matching loss \citep{koehler2022statisticalefficiencyscorematching}. 

This idea is illustrated in \cref{fig:scores}. The top row plots the score functions for two Gaussian mixture distributions that are identical except for the mixture weight $\gamma$ and shows that they are nearly identical except in a region with very low probability mass. We are thus unlikely to have sufficient training samples for the loss to reliably indicate the value of $\gamma$ that best represents the training data. However, for larger values of $t$ in the diffusion process ($t=0.25$ or $t=0.50$ in the figure), the corresponding distributions are less well separated, and the score functions are clearly distinguishable over a higher-probability region. Since training diffusion models uses the score across multiple $t \in [0,1]$, learned models hone in on the correct value of $\gamma$.

\begin{figure}[ht]
    \centering\includegraphics[width=0.7\linewidth]{images/figure_zero.pdf}
    \caption{
    Score sensitivity to mixture weights at different diffusion times. At $t=0$, the scores for $\gamma=0.8$ and $\gamma=0.5$ are nearly identical except on a low-probability region. At $t>0$, the mixture components overlap more, and the region where the scores differ has a substantially higher probability.
    % Score sensitivity to mixture weights at different stages of the diffusion process. The original mixture distribution in the top row corresponds to mixture weight $\gamma^*=0.8$; if we look at the score function for this mixture compared to the score function of a distribution with the same components but mixture weight $\gamma = 0.5$, the scores are nearly identical except for a very low probability subset of the domain. But for $t>0$, the corresponding distributions are less separated, and the probability mass of the part of the domain where the scores are distinct is much higher.
    }
    \label{fig:scores}
\end{figure}

To quantify this phenomenon, we measure the sensitivity of the diffusion score matching loss $\ldsm$  \eqref{eq:diffScoreMatchingLoss} to changes in parameters of interest. 
Specifically, let $\{\pdist{\gamma}\}_{\gamma \in \bm{\Gamma}}$ be a set of parameterized distributions. Let $\pdist{\gamma^*}$ be a target distribution with target parameter $\gamma^*$. %Let $\ldsm$ denote the diffusion score matching loss (annealed over multiple $t$, , %where $\gamma^*$ is the true parameter of interest, 
%then 
The \textit{diffusion score sensitivity index (DSSI)} of the target parameter $\gamma^*$ %w.r.t a target parameter $\pdist{\gamma^*}$
is defined as,
\begin{equation}\label{eq: DSSI}
    L(\gamma^*) = \inf_{\gamma \in \bm{\Gamma}}~ \frac{\ldsm(\pdist{\gamma},\pdist{\gamma^*})}{(\gamma-\gamma^*)^2}
\end{equation}
 
We note that both $\ldsm$ and $L(\gamma^*)$ depend on our choice of noise schedule.
If the target data distribution has mixture weight $\gamma^*$, and we learn a distribution with mixture parameter $\gamma$, then the squared error of mixture weight estimation $\gamma$ is bounded by the diffusion score matching loss divided by the DSSI. Thus, when the DSSI is small, our learned distribution may be highly misrepresentative of $\gamma^*$ even when the diffusion score matching loss is low. In contrast, when $L(\gamma^*)$ is large, we know that a small diffusion score matching loss indicates low error in the parameter $\gamma$. Our main contributions are as follows:
\begin{itemize}
    \item For general bimodal mixtures, we prove a strict minimax separation in mixture-weight recovery between estimators that use only the target-score data and estimators with access to score data at an intermediate, mode-overlapping noise level. (Theorem \ref{thm:sensitivity})
    \item For Gaussian mixtures, we prove a lower bound on the DSSI that holds in any dimension; thus, parameter estimation error is controlled by the DSM loss (Theorem \ref{thm: GMM Sensitivity Emerges}).
    
    \item We further empirically evaluate the DSSI on a real-world dataset. These experiments illustrate how design decisions relating to diffusion model training and inference-time sampling schedules can affect the DSSI and the fidelity of generated samples to the training data distribution.
\end{itemize}
\section{Preliminaries: score-based diffusion models}
\label{sec:prelim}
Score-based diffusion models sample from a target distribution on $\mathbb{R}^d$ by tracing out a one-parameter family of distributions $\{p_t\}_{t \in [0,1]}$ indexed by a diffusion time $t$, with $p_0$ being the target data distribution and $p_1 \approx \mathcal{N}(\mathbf{0}, \mathbf{I})$ the standard normal distribution. A \emph{forward process} progressively transforms samples from $p_0$ into samples from $p_1$ along this family. A \emph{reverse process} inverts the densities along the forward process, starting from $\mathcal{N}(0,\mathbf{I})$, using learned approximations of the scores of the intermediate marginals $p_t$.
Thus, score-based diffusions have 3 components: (a) the forward diffusion process and its corresponding reverse process, (b) a training objective for learning the marginal score functions $\nabla \log p_t$ of the intermediate distributions for $t\in [0,1]$,  and (c) a sampling procedure used to approximately generate samples from the target data distribution at inference time. In this section, we provide a brief description of each component.



\paragraph{Noising process.} 
Given a data distribution $p_0$ on $\mathbb{R}^d$, the noising process we consider is the {variance-preserving forward SDE} (VP-SDE)  proposed in \citet{song2021scorebasedgenerativemodelingstochastic}. Give a noise schedule $\set{\beta_t}$, this takes the form 
\begin{equation}
    \mathrm{d}{\vx}_t 
    \;=\; 
    -\tfrac{1}{2} \beta_t \, {\vx}_t\, \mathrm{d}t \;+\; \sqrt{\beta_t}\,\mathrm{d}\vect{w}_t,
    \qquad {\vx}_0 \sim p_0,\quad t \in [0,1],
    \label{eq:forward-sde}
\end{equation}
where $\vect{w}_t$ is a standard $d$-dimensional Wiener process. We denote by $p_t$ the marginal distribution of ${\vx}_t$. The forward conditional law of ${\vx}_t$ given ${\vx}_0$ under \eqref{eq:forward-sde} is Gaussian, 
\begin{equation}
    p_t({\vx}_t \mid {\vx}_0) 
    \;=\; 
    \mathcal{N}\!\left(
    {\vx}_t;\; \sqrt{\alpha_t}\,{\vx}_0,\; (1-\alpha_t)\,\vect{I}
    \right),
   \label{eq:forward-kernel}
\end{equation}
with conditional mean $\sqrt{\alpha_t} {\vx}_0$ and conditional covariance $(1-\alpha_t) \vect{I}$, where $\alpha_t=e^{-\int_0^t\beta_s \mathrm{d}s}$.\footnote{As a slight abuse of terminology, the sequences $\set{\alpha_t}$ and $\set{\beta_t}$ will interchangably be referred to as the ``noise schedule''.} 

\paragraph{Denoising process.}
\citet{ANDERSON1982313} established that the time-reversal of \eqref{eq:forward-sde} is itself an SDE:
% , drive by the marginal score $\score[t]{\gamma}$:
\begin{equation}
    \mathrm{d}{\vx}_t 
    \;=\; 
    \left[-\tfrac{1}{2}{\beta_t}\,{\vx}_t \;-\; {\beta_t}\,\nabla_{{\vx}} \log p_t({\vx}_t)\right]
    \mathrm{d}t 
    \;+\; 
    \sqrt{{\beta_t}}\,\mathrm{d}\bar{\vect{w}}_t,
    \qquad {\vx}_1 \sim p_1,
    \label{eq:reverse-sde}
\end{equation}
where $\bar{\vect{w}}_t$ is the reverse-time Wiener process. 
% As per \eqref{eq:reverse-sde}, s
Starting from ${\vx}_1 \sim p_1$ and integrating the reverse SDE with the exact score function from time $t=1$ to $t=0$ produces samples ${\vx}_0 \sim p_0$.  In practice, this reverse SDE is simulated starting from ${\vx}_1 \sim \mathcal{N}(0,\mathbf{I}) \approx p_1$ using the learned approximation $\vs_{\btheta}: \mathbb{R}^d \times [0,1] \rightarrow \mathbb{R}^d$ with parameters $\btheta  \in \boldsymbol{\Theta}$ trained such that $\vs_{\btheta}({\vx},t) \approx \nabla \log p_t({\vx})$ simultaneously for all noise levels to produce samples approximately distributed as the target $p_0$. % $t\in [0,1]$. 
% For a review of the objectives used to learn score networks and brief remarks on the significance of certain design choices, see \Cref{Sec: SM Loss}. 



For a given trained score model $\{\vs_{\btheta }(\cdot, t)\}_{t\in [0,1]}$ (or the true score function $\{\nabla \log p_t\}_{t\in [0,1]}$), generation of samples is an inference task: one may simulate the reverse SDE in \eqref{eq:reverse-sde} with any choice of consistent numerical integration schemes \citep{karras2022elucidating}.
The overall process has 
% Hence, the quality of the generated samples can be attributed to 
two distinct sources of error: the \emph{approximation error} incurred in learning the score
% , inherited from \eqref{eq:dsm-loss}, 
and the \emph{sampling error} incurred by the choice of sampling algorithm (e.g., discretization error, finite timesteps, etc.). 

 


\section{Minimax bounds for parameter estimation from diffusion models}\label{sec: parameterized mixture families}
In this section, we show that when the score along the noising process exhibits sensitivity to $\gamma$, the information-theoretic barrier that disallows its recovery from the target score alone is removed. We prove that when the score is sensitive at some noise level, we obtain a lower parameter estimation error when compared to only using the target score. 
% as a consequence, we incur lower parameter estimation error. 
While stated explicitly for mixture weight as the parameter, the results in this section apply to any parameter. We consider only bimodal mixtures for simplicity, but make no distributional assumptions about the components. 
\paragraph{Mixture Family.}
Let $\pdist{0}$ and $\pdist{1}$ be two fixed probability distributions on $\mathbb{R}^d$. A probability density $\pdist{\gamma}, \gamma \in \bm{\Gamma} \subseteq (0,1),$ is a member of a family of two-component mixtures if it is of the form $\pdist{\gamma}\coloneqq (1-\gamma)\: \pdist{0} + \gamma \:\pdist{1}$.
The mixture parameter $\gamma$ controls the relative weights of the two component distributions, $\pdist{0}$ and $\pdist{1}$. We now consider the evolution of a mixture, $\pdist{\gamma},$ under a noising process or forward process in a score-based diffusion model. 
Specifically, when $p_0 \equiv \pdist{\gamma},$ we denote by $\pdist[t]{\gamma}$ the marginals of $\vx_t$ following the process \eqref{eq:forward-sde}. Let $\score{\gamma}$ denote the {score} of the density $\pdist{\gamma}$, and similar for marginal densities of $\vx_t$:
% which is a vector field defined pointwise as 
\begin{equation}\label{eq:score-and-noisyscore}
  \score{\gamma} := \nabla \log \pdist{\gamma} \qquad \text{and} \qquad \score[t]{\gamma} := \nabla \log \pdist[t]{\gamma}.  
\end{equation}
%$$$$
The relative weights of the mixture are preserved under the forward process: the noisy marginal $\pdist[t]{\gamma}$ is itself a $\gamma$-mixture of the noisy component marginals $\pdist[t]{0}$ and $\pdist[t]{1}$ (see \Cref{lemma:mixture-closure}). 
% Using an analogous notation, w
% Furthermore, we write $\score[t]{\gamma} := \nabla \log \pdist[t]{\gamma}$ for the score corresponding to marginal densities of $\vx_t$ satisfying \eqref{eq:forward-sde}. 
% \becca{careful here because here we're using $w$ for weights but in \eqref{eq:forward-sde} we use $\vect{w}$ for Wiener process. }
As a consequence the score of the noisy mixture density $\pdist[t]{\gamma}$ can be decomposed as follows,
\begin{align}
\score[t]{\gamma}({\vx})
    = w_t^{(\gamma,0)}({\vx})\,\score[t]{0}({\vx}) \;+\; w_t^{(\gamma,1)}({\vx})\,\score[t]{1}({\vx})
    \label{eq:mix score}
\end{align}
% Taking derivatives \becca{wrt what? do you mean computing the score of \eqref{eq:mixture-closure}?}
% of \eqref{eq:mixture-closure}, the noised mixture score $\score[t]{\gamma}$ is a pointwise convex combination of the component scores: i.e., $\score[t]{\gamma}({\vx})
%     = w_t^{(0)}({\vx})\,\score[t]{0}({\vx}) \;+\; w_t^{(1)}({\vx})\,\score[t]{1}({\vx})$
with component weights $w_t^{(\gamma,0)}({\vx}) = (1-\gamma)\,\pdist[t]{0}({\vx})/\pdist[t]{\gamma}({\vx})$ and $w_t^{(\gamma,1)}({\vx}) = \gamma\,\pdist[t]{1}({\vx})/\pdist[t]{\gamma}({\vx})$.
% satisfy $w_t^{(0)}({\vx}) + w_t^{(1)}({\vx}) = 1$. 
This decomposition enables a simplified theoretical analysis, as we see in the next section. 

\paragraph{Classical score matching loss/relative Fisher information.} The classical score matching (SM) loss measures the relative Fisher information between distributions \citep{10.5555/1046920.1088696}. For a pair of mixtures, % corresponding to mixture weights $\gamma, \gamma'$  :
\begin{align}
    \lsm(\pdist{\gamma'},\pdist{\gamma}) := \mathbb{E}_{\vx \sim \pdist{\gamma}} \|\score{\gamma'}(\vx)-\score{\gamma}(\vx) \|^2.
\end{align}
The relative Fisher information between a pair of mixture distributions with identical components but different mixture weights is small when the components are well-separated. In \cref{fig:scores}, the scores of Gaussian mixtures may appear nearly identical outside a low-probability region, an effect strengthened with increasing separation between mixture components as shown in \cref{fig: different loss functions}.   %also shows that the relative Fisher information, or equivalently $\ell_\mathrm{SM}$ varies inversely with mode separation. 
% for 2D Gaussian mixtures.
% Insensitivity of the loss with respect to parameters at high separation implies the resulting generated samples will also be insensitive. 
Any sampling algorithm that relies solely on the score, such as Langevin sampling, exhibits near-identical behavior across distinct mixture distributions; consequently, the generated samples may not accurately reflect the target mixture weights.  %from the two mixtures and cannot distinguish them.
%Thus samples generated using algorithms that rely solely on the score of a target distribution $\pdist{\gamma}$,  inherit score of a  Fisher information, such as Langevin sampling, are insufficient for distinguishing between nearby mixture distributions. % producing samples even with an accurate score can lead to samples with the incorrect mixture weight. 

\paragraph{Diffusion score matching loss for mixtures.} 
%The two parts of a generative modeling algorithm that influence parameter identifiability are the loss function and the sampling process. 
The diffusion score matching (DSM) loss anneals the score matching loss across the noising process \eqref{eq:forward-sde} \citep{song2019generative}. For a pair of mixtures, % distributions  $\pdist{\gamma}$ and $\pdist{\gamma'}$, %measures the average score matching loss over the corresponding noising processes,
\begin{align}
    \ldsm( \pdist{\gamma'},\pdist{\gamma}) &:=  \mathbb{E}_{\substack{t \sim \mathrm{Unif}(0,1)\\\vx\sim \pdist[t]{\gamma}}}[\lambda_t \|\score[t]{\gamma'}(\vx) - \score[t]{\gamma}(\vx)\|^2] 
    = \mathbb{E}_{t \sim \mathrm{Unif}(0,1)}[\lambda_t \: \lsm(\pdist[t]{\gamma'},\pdist[t]{\gamma})].  \label{eq:diffScoreMatchingLoss}
\end{align}
Here, $\lambda_t$ is a weighting factor set to $\lambda_t = 1-\alpha_t$ following \citet{song2021scorebasedgenerativemodelingstochastic}. The DSM loss measures the relative Fisher information over the entire sequence of noised densities generated by the process \eqref{eq:forward-sde}. Recall that the marginals $p_t$'s and scores $\vs_t$'s depend on the noise schedule $\{\alpha_t\}_t$, as shown in \eqref{eq:forward-kernel} and \eqref{eq:score-and-noisyscore}, and this way, the diffusion score matching loss $\ldsm$ is dependent on the choice of noise schedule. %, even though this is not explicit in our notation.
In practice, %he true score function $\score[t]{\gamma}$ is unknown, and thus calculating 
we may train a neural network $\vs_{\btheta}$ to approximate the score function. With some abuse of notation, we denote $\pdist{\btheta}$ the distribution\footnote{The superscript here indicates the trainable parameters $\btheta$ of the learned network.} of samples generated using the score functions $\{\vs_{\btheta}(\cdot, t)\}_{t\in (0,1)}$ for a fixed noise schedule as per \eqref{eq:reverse-sde}. Since the true score function $\score[t]{\gamma}$ and indeed the target distribution $\pdist{\gamma}$ are unknown,  calculating the diffusion score matching loss %$\ldsm(\vect \vs_{\btheta}, \pdist[t]{\gamma})$ 
is intractable. %Due to Tweedie's formula, we may reformulate \eqref{eq:diffScoreMatchingLoss} in terms of Using 
Instead, we minimize the objective 
\begin{equation}\label{eqn: neural network loss}
\lnn(\pdist{\btheta}, \pdist{\gamma}):=\mathbb{E}_{\substack{t \sim \mathrm{Unif}(0,1)\\\vx_0\sim \pdist{\gamma}, \vx_t\sim \pdist[t]{\gamma}(\cdot | \vx_0)}}[\lambda_t \|\vect \vs_{\btheta}(\vx_t, t) - \nabla \log \pdist[t]{\gamma}(\vx_t|\vx_0)\|^2].
\end{equation}



As per \Cref{{lem:nn-dsm-ordering}}, there exists a constant $C\geq 0$, depending only on $\pdist{\gamma}$, such that $\lnn(\pdist{\btheta}, \pdist{\gamma})=\ldsm(\pdist{\btheta}, \pdist[t]{\gamma}) + C$. % for a positive constant $C$, which is only dependant on the data distribution $\pdist{\gamma}$. 
Thus, %We caution that using $\lnn<\delta$ as a proxy for $\ldsm<\delta$ is conservative, as the two functions differ by a constant depending on $\pdist{\gamma^*}$.  
$\ldsm(\pdist{\btheta}, \pdist{\gamma}) \leq \lnn(\pdist{\btheta}, \pdist{\gamma})$.

\paragraph{Minimax theory for mixture weight recovery.} From classical minimax theory (see e.g., \citep{geer2000empirical} for a textbook exposition), we know that the error in any estimator of a parameter is lower-bounded by the informativeness of the data distribution about the parameter. Now, the data distribution loses information about the parameter in the noising process. How, then, would annealing benefit  arameter recovery, if the original mixture is well-separated and hence insensitive to parameter changes? The answer lies in the fact that diffusion models learn to approximate the \emph{scores} at various noise levels. By minimizing $\ldsm,$ diffusion models obtain information about the parameter through the scores even if such information is not present in the samples.  Our main result in this section shows the existence of a noise scale at which the score-data law is sensitive, provably enabling parameter recovery rates that are unattainable from target-score data alone.


\begin{definition}[Score oracle and score-restricted estimators] We consider a Gaussian White Noise (GWN) oracle from nonparametric regression \cite{tsybakov2008nonparametric} of the form: $\hat{\mathbf s}_{t}^{(\gamma)} = \score[t]{\gamma}\:  + \sigma_t \: W,$ where $\sigma_t > 0$ is a scalar and $W$ is an isonormal field indexed by $L^2(\pdist[t]{\gamma})$. The statistical experiment is to estimate $\gamma$ from the observation, $\hat{\mathbf s}^{(\gamma)}_t.$ Let $\qdist[t]{\gamma}$ denote the distribution of the observation, which is a Gaussian measure with mean $\score[t]{\gamma}.$ We define the set, $\Gamma_{\mathrm{score},t},$ of all measurable functions of the GWN observations to $(0,1)$ to be the class of ``score-restricted estimators''.  Further, we use $ \mathbb{E}_{\gamma^*,t}$ as a shorthand for expectations with respect to the score oracle model, $\qdist[t]{\gamma}.$ See also Remark \ref{remark:score-restricted-estimators} for more details on the estimator classes.
\label{def:scoreEstimator}
\end{definition}

We call $\Gamma_{\mathrm{score},t}$ \emph{score-restricted estimators} since they access the score vector field (corrupt by Gaussian noise) and not the target data samples. In other words, our class of estimators $\Gamma_{\mathrm{score},t}$ models parameter recovery achievable from having oracle access to a corrupt score vector field at time $t,$ reflecting the denoising/generation phase in diffusion models after scores have been trained. 
Since the observation model, $\qdist[t]{\gamma},$ is Gaussian, we can compute, fixing the indexing Hilbert space for $W$ to be $L^2(\pdist[t]{\gamma}),$ that $\mathrm{KL}(\qdist[t]{\gamma}||\qdist[t]{\gamma'}) = \lsm(\pdist[t]{\gamma'}, \pdist[t]{\gamma})/(2\sigma_t^2).$ See Remark \ref{remark:cameronMartin} for this computation. 

\begin{assumption}[Stable score-matching at time $t$] We assume that at some time $t > 0,$ the score observation law, $\qdist[t]{(\cdot)},$ satisfies $ (\gamma - \gamma')^2 \leq A_t^2 \: \mathrm{KL}(\qdist[t]{\gamma}||\qdist[t]{\gamma'})$ for some $A_t >0$ and  all $\gamma, \gamma' \in (0,1).$  
\label{assumption:continuity}
\end{assumption}

Since $\mathrm{KL}(\qdist[t]{\gamma}||\qdist[t]{\gamma'}) = \lsm(\pdist[t]{\gamma'}, \pdist[t]{\gamma})/(2\sigma_t^2),$ the above assumption implies that when the score matching loss is bounded above by some $C,$ then, the squared difference in the corresponding parameters, $(\gamma - \gamma')^2 \leq  \:A_t^2\: C / (2\sigma_t^2).$ 
% \becca{I think this might be a mistake:
% \begin{align}
%     (\gamma-\gamma')^2 &\le A_t^2 {\rm KL}(q_t^{(\gamma)}\|q_t^{(\gamma')})\\
%     &= A_t^2 \lsm(p_t^{(\gamma)},p_t^{(\gamma')})/(2\sigma_t^2)\\
%     &\le A_t^2 C/(2\sigma_t^2)
% \end{align}
% }
Intuitively, this means that, for an 
% $O(1/\sigma_t^2)$ 
% \becca{
$O(\sigma_t^2)$
constant $A_t,$ the score observations predict distances on the space of weight parameters.    Our next definition allows us to make this setting rigorous by comparing score observations at  time $t$ to a different behavior of the scores at time $0:$ the target score observations, $\hat{\mathbf s}^{(\gamma)},$ are insensitive in law to changes in $\gamma$ and therefore, do not predict distances on parameter space. In the next section, note that we obtain explicit estimates on $A_t$ for Gaussian mixtures. 

\begin{definition}[$(\epsilon,\delta)$-diffusion sensitivity]
\label{definition:diffusionSensitivity} We say that a mixture family $\{\pdist{\gamma}\}$ exhibits an $\epsilon$-diffusion sensitivity if for every $\gamma^* \in (0,1),$ there exist $\epsilon$-packing sets $S(\gamma^*) \subset (0,1)$ containing $N(\gamma^*,\epsilon) \geq 16$ elements with 
$\mathrm{KL}(\qdist{\gamma}||\qdist{\gamma'}) = \lsm(\pdist{\gamma'}, \pdist{\gamma})/(2\sigma_0^2) \leq \delta(\epsilon,N) \leq \epsilon^2/4,$ for all $\gamma, \gamma' \in S(\gamma^*)$ satisfying $|\gamma - \gamma'| \geq  4\sqrt{2}\: A_t \epsilon$. 
\begin{comment}\becca{the role of $t$ is unclear here. Up until this point, the definition described a property of a mixture family. Which $t$ do we mean? Do we just mean that we have $(\epsilon,\delta)$-sensitivity if there exists some $t \in [0,1]$ for which this is true? Or maybe you mean we have $(\epsilon,\delta)$-sensitivity if the KL bound in this definition holds for all $t$ for which \cref{assumption:continuity} holds?}
This holds for a specific $t$ that satisfies the assumption above.
\end{comment}
Note that $A_t$ here is as defined in Assumption \ref{assumption:continuity}.\end{definition}

\begin{definition}[$\epsilon$-diffusion sensitivity]
\label{definition:diffusionSensitivity} We say that a mixture family $\{\pdist{\gamma}\}$ exhibits $\epsilon$-diffusion sensitivity if for every $\gamma^* \in (0,1),$ there exist $\epsilon$-packing sets $S(\gamma^*) \subset (0,1)$ containing $N(\gamma^*,\epsilon) \geq 16$ elements with 
 $\mathrm{KL}(\qdist{\gamma}||\qdist{\gamma'}) = \lsm(\pdist{\gamma'}, \pdist{\gamma})/(2\sigma_0^2) \leq \epsilon^2/4,$ for all $\gamma, \gamma' \in S(\gamma^*)$ satisfying $|\gamma - \gamma'| \geq  4\sqrt{2}\: A_t \epsilon$. Note that $A_t$ here is as defined in Assumption \ref{assumption:continuity}.\end{definition}
 




\begin{theorem}
\label{thm:sensitivity}
Fix $t>0$ satifying Assumption \ref{assumption:continuity}. Let $\epsilon_{t} > 0$ be a covering radius that solves $\epsilon_{t}:= \sigma_t\: \sqrt{V(\epsilon_{t})},$ where $V(\epsilon_{t})$ is the $\epsilon_{t}$-covering entropy\footnote{Refer \Cref{def:cover-entropy} for an explicit statement} with respect to square root KL divergence of the class of Gaussian measures $\{\qdist[t]{\gamma}\}.$ 
     Suppose the bimodal mixture family satisfies \Cref{assumption:continuity} and exhibits an $\epsilon$-diffusion sensitivity with $\epsilon = \epsilon_{t}.$
    Then, we have the following minimax upper bound on the score-restricted estimators at time $t$, 
    \begin{align}
        \min_{\hat{\gamma} \in \Gamma_{\mathrm{score},t}}\max_{\gamma^*}~
        \mathbb{E}_{\gamma^*,t} |\gamma^* - \hat{\gamma}|^2 \leq 2\:A_t^2\: \epsilon_{t}^2\: .
        \label{eq:upperBound}
    \end{align} 
Simultaneously, if the packing set cardinality $N(\gamma^*, \epsilon_{t}) \geq e^{\sigma_t^2\:V(\epsilon_{t})},$ for all $\gamma^*,$ it  holds that the minimax lower bound achievable for target score-restricted estimators is the following:
    \begin{align}
\label{eq:lowerBound}
\min_{\hat{\gamma} \in \Gamma_{\mathrm{score},0}}\max_{\gamma^*}~ \mathbb{E}_{\gamma^*,0} |\gamma^* - \hat{\gamma}|^2 \geq 4\: A_t^2\: \epsilon_{t}^2.
    \end{align}
\end{theorem}


% \begin{theorem}
% \label{thm:sensitivity}
%     Let $\epsilon_{m,t} > 0$ be a covering radius that solves $\epsilon_{m,t}:= \sqrt{V(\epsilon_{m,t})/m},$ where $V(\epsilon_{m,t})$ is the $\epsilon_{m,t}$-covering entropy\footnote{Refer \Cref{def:cover-entropy} for an explicit statement} with respect to square root KL divergence of the class of densities $\{\qdist[t]{\gamma}\}.$ 
%      Suppose the bimodal mixture family satisfies \Cref{assumption:continuity} and \Cref{definition:diffusionSensitivity} at time $t$ with $(\epsilon, \delta) = (\epsilon_{m,t}, V(\epsilon_{m,t})/4m).$
%     Then, we have the following minimax upper bound on the score restricted estimators using $m$ samples of score data from $\{\qdist[t]{\gamma}\}$ at time $t$:
%     \begin{align}
%         \min_{\hat{\gamma} \in \Gamma^{(m)}_{\mathrm{score},t}}\max_{\gamma^*}~
%         \mathbb{E}_{\gamma^*,t} |\gamma^* - \hat{\gamma}| \leq \sqrt{2}\:A_t\: \epsilon_{m,t}\: .
%         \label{eq:upperBound}
%     \end{align} 
% Simultaneously, it holds that the minimax lower bound achievable for target score-restricted estimators is the following:
%     \begin{align}
% \label{eq:lowerBound}
% \min_{\hat{\gamma} \in \Gamma^{(m)}_{\mathrm{score},0}}\max_{\gamma^*}~ \mathbb{E}_{\gamma^*,0} |\gamma^* - \hat{\gamma}| \geq 2\: A_t\: \epsilon_{m,t}
%     \end{align}
% \end{theorem}

\paragraph{Proof.} The upper bound in \eqref{eq:upperBound} is a direct application of Theorem 2 from \citep{yangBarron}. Then, using a generalized Fano's method (Lemma 3 of \citep{binyu}) and an appropriate choice of $\delta$ as a function of $\epsilon$ in \cref{definition:diffusionSensitivity} that limits the packing entropy of the cover of $(0,1)$ with cardinality $e^{V(\epsilon_{t})},$ we obtain the desired lower bound in \eqref{eq:lowerBound}. We defer a more detailed proof to  Appendix \ref{sec:lowerUpperBounds}. \qed

Theorem \ref{thm:sensitivity} elucidates the conditions under which a diffusion model learns the mixture weights accurately when target score estimation followed by Langevin sampling does not. The above analysis 
applies to any bimodal mixture density and to any scalar parameter in place of the mixture weight. 
The main mechanism that enables parameter recovery is that the score field at some time $t$
changes more rapidly with parameter changes. In contrast, the target score at time 0 may not vary, creating an information-theoretic barrier that limits the accuracy of any
estimator. That is, no estimator can distinguish two close values of $\gamma$ using a finite number of target
score samples. 


\begin{remark}
A key contribution of our work is to explain the surprising phenomenon in which a diffusion model can accurately recover a parameter even when the target score remains insensitive to small perturbations of the parameter. In response, we defined the notion of DSSI, \eqref{eq: DSSI}, which measures the sensitivity of the scores along the entire noising process, as opposed to only the target score, to infinitesimal changes in the parameter. The most interesting scenario that is directly relevant to explaining ``accurate parameter recovery despite target score insensitivity'' is one where an intermediate score is ``much more'' sensitive to the parameter than the score of the target distribution. Our definition of the $\epsilon$-sensitivity is designed to capture this scenario and further allows us to fully explain the information-theoretic mechanism underlying accurate recovery (Theorem \ref{thm:sensitivity}). 
\begin{comment}     
    Note that $(\epsilon,\delta)$-sensitivity is a measure of sensitivity occurring at later times than $t=0$, while the DSSI $L(\gamma)$ only measures average sensitivity over the full noising process. However, under weak assumptions, for a distribution where the DSSI $L(\gamma^*)$ is large, a local version of \cref{assumption:continuity} holds for a small value of $A_t$. The conditions for exhibiting an $(\epsilon,\delta)$-sensitivity are easier to satisfy when $A_t$ is small, so this suggests that families of distributions such that $L(\gamma)$ is large for all $\gamma$ are more likely to exhibit an $(\epsilon,\delta)$-sensitivity. \end{comment}
\end{remark}


% \todo[inline]{Even though this is an expectation over time, we know that there exists a time $t$ where a large sensitivity emerges if DSSI is large. If additionally $A_t$ is small relative to $A_0$, then $(\epsilon,\delta)$-sensitivity also holds.}

% \todo[inline]{Why we care: The $(\epsilon,\delta)$-sensitivity is a more fine-grained notion in time, and is motivated to explain that a sensitivity could emerge during the noising process, even if it is negligible at time $t=0$. It's here to justify that the diffusion model can develop sensitivity in practice.}
% \todo[inline]{This is a mechanistic understanding of when we get more information using the forward process compared to just the data distribution. The DSSI, on the other hand, is more computable.}

% \todo[inline]{Change uses of Assumption with Definition in the text}

We remark that \citet{yangBarron} constructs a Bayes density estimator that belongs to the mixture class and achieves the upper bound \eqref{eq:upperBound}. Thus, we can conclude from Theorem \ref{thm:sensitivity} that using score observations at time $t$ leads to provably better parameter recovery.
 

%That is, no estimator can distinguish two close values of $\gamma$ using a finite number of target score samples. 

\textbf{Resolving data-processing inequality.} Saying a family of distributions exhibits an $(\epsilon,\delta)$-sensitivity essentially means that more information about the parameter is created in the noising process, which, at first glance, appears to violate the data processing inequality. The law of noised samples, $\pdist[t]{\gamma},$ cannot be more informative about $\gamma$ at any time $t$ than $\pdist{\gamma}$ (the target mixture), due to the data processing inequality. More precisely, since $\pdist[t]{\gamma}$ results from a $\gamma$-independent convolution operation on $\pdist[0]{\gamma},$ the data processing inequality gives: 
$\mathrm{KL}(\pdist[t]{\gamma'}||\pdist[t]{\gamma}) \leq \mathrm{KL}(\pdist{\gamma'}||\pdist{\gamma}).$ But, our key idea here is that diffusion models use the score observation fields, $\hat{\mathbf s}^{(\gamma)}_t,$ to recover $\gamma$, even when $\hat{\mathbf s}^{(\gamma)}$ can not separate $\gamma$ values sufficiently. Specifically, under the assumptions of Theorem \ref{thm:sensitivity} on the mixture family, the score observations can be more informative about $\gamma$ at time $t$ than at time 0, i.e., $\mathrm{KL}(\qdist[t]{\gamma'}||\qdist[t]{\gamma}) \geq \mathrm{KL}(\qdist{\gamma'}||\qdist{\gamma}),$ without violating data processing. In other words, \cref{definition:diffusionSensitivity} defines a structural property of mixture distributions: the family, $\{\qdist[t]{\gamma}\}$ at intermediate $t$ can be more separated in $\gamma$ than at $t=0$ when mixture components do not overlap.
 In Lemma \ref{remark:dataProcessing}, we prove why this counterintuitive inequality does not violate data processing and why no estimator in $\Gamma_0$ can reconstruct the score observation model that $\Gamma_t$ has access to. 
% \becca{I would like to say something perhaps a little more directly here, along the lines of: Let ${\cal K}$ correspond to any operator that can be applied to $q^{(\gamma)}$ for any $\gamma$. By the data processing inequality, we have
% \begin{align}
%     {\rm KL}({\cal K} q_0^{(\gamma)}\|{\cal K} q_0^{(\gamma')}) &\le {\rm KL}(q_0^{(\gamma)}\| q_0^{(\gamma')}).
% \end{align}
% Simultaneously, under the assumption of $\epsilon$-sensitivity, we know 
% \begin{align}
% {\rm KL}(q_0^{(\gamma)}\| q_0^{(\gamma')}) \le {\rm KL}(q_t^{(\gamma)}\| q_t^{(\gamma')}) / 128.
% \end{align}
% To see this, note
% \begin{align}
% {\rm KL}(q_0^{(\gamma)}\| q_0^{(\gamma')}) \le & \frac{\epsilon^2}{4} 
% & \text{by $\epsilon$-sensitivity}
% \\
% \le & \frac{(\gamma-\gamma')^2}{4(4\sqrt{2})^2A_t^2}
% & \text{by $\epsilon$-sensitivity}
% \\
% = & \frac{(\gamma-\gamma')^2}{128A_t^2}\\
% \le & \frac{A_t^2 {\rm KL}(q_t^{(\gamma)}\| q_t^{(\gamma')})}{128 A_t^2}
% & \text{by Assumption 3.2}
% \\
% = &\frac{{\rm KL}(q_t^{(\gamma)}\| q_t^{(\gamma')})}{128}.
% \end{align}
% Putting the pieces together, we have
% \begin{align}
%     {\rm KL}({\cal K} q_0^{(\gamma)}\|{\cal K} q_0^{(\gamma')}) &\le {\rm KL}(q_0^{(\gamma)}\| q_0^{(\gamma')}) \le \frac{{\rm KL}(q_t^{(\gamma)}\| q_t^{(\gamma')})}{128}.
% \end{align}
% When Nisha and I spoke, she suggested $q_t^{(\gamma)}$ could not be computed from $q_0^{(\gamma)}$ without violating this chain of KL inequalities, and this is why we can think of estimators that only use $q_0^{(\gamma)}$ as separate from the class of estimators that use $q_t^{(\gamma)}$, and so those two classes of estimators could have different rates. But because of the $1/128$ factor, I'm worried that in fact ${\cal K}$ \textit{could} map to $q_t^{(\gamma)}$, and so it's not totally clear to me that the class of estimators that only use $q_0^{(\gamma)}$ should have different rates.
% }




\subsection{Effect of modified sensitivity to mixture weight}\label{sec: GMM Sens Exp}

While Theorem \ref{thm:sensitivity} provides a mechanism for parameter recovery, we now discuss experimental evidence for the same mechanism but using the practically computable diffusion score sensitivity index \eqref{eq: DSSI}, rather than the information-theoretic quantities used in \cref{definition:diffusionSensitivity}. Our numerical results show that the score sensitivity index quantifies how accurately a diffusion model captures the mixture weight of a Gaussian mixture model. We generate samples using the true analytical score function $\score[t]{\gamma^*}(\vect x)$ and a sequence of modified noise schedules with a coarse time discretization ($T=100$), designed to artificially reduce the diffusion score sensitivity index $L(\gamma^*)$ (for details, see \Cref{sec: Reparam}). Recall that $L(\gamma^*)$ and $\ldsm$ depend on the choice of noise schedule $\{\alpha_t\}_t$. %, as discussed below \eqref{eq:diffScoreMatchingLoss}.

We modify both linear and squared-cosine noise schedules and numerically estimate $L(\gamma^*)$ for these schedules.  \Cref{fig: gmm est gamma vs sensitivity index} shows the estimated mixture weight $\hat\gamma$ (calculated using Expectation-Maximization) plotted against $L(\gamma^*)$. We see that as the sensitivity index $L(\gamma^*)$ decreases, both the bias and variance in $\hat\gamma$ increase. This demonstrates that noise schedules with lower $L(\gamma^*)$ yield mixture-weight estimates that are less robust to discretization errors, even with a perfect score estimator.

On the other hand, the means and covariances of the modes are uniformly well-approximated across the different noise schedules. This shows that errors in estimating the mixture weight can occur without noticeable artifacts in the geometry of individual modes. Modern accelerated diffusion sampling schemes (e.g., \citep{song2021denoising}) leverage coarse time discretizations to reduce the computational cost of sampling while maintaining comparable sample fidelity. Our results suggest that such sampling approaches may be vulnerable to misrepresenting parameters of interest in the data distribution while appearing to generate accurate samples.

\begin{figure}
    \centering
    \begin{subfigure}[htbp]{0.24\linewidth}
        \includegraphics[width=\columnwidth]{images/figure_two_a.pdf}
        \caption{}
        \label{fig: gmm lin heatmap}
    \end{subfigure}
    \hfill
    \begin{subfigure}[htbp]{0.24\linewidth}
        \includegraphics[width=\columnwidth]{images/figure_two_b.pdf}
        \caption{}
        \label{fig: gmm score sensitivity}
    \end{subfigure}
    \hfill    
    \begin{subfigure}[htbp]{0.24\linewidth}
        \includegraphics[width=\columnwidth]{images/figure_two_c.pdf}
        \caption{}
        \label{fig: gmm est gamma vs sensitivity index}
    \end{subfigure}
    \hfill    
    \begin{subfigure}[htbp]{0.24\linewidth}
        \includegraphics[width=\columnwidth]{images/figure_one_e.pdf}
        \caption{}
        \label{fig: different loss functions}
    \end{subfigure}
    
    \caption{
    Score sensitivity for Gaussian mixture models in dimension $d=10$. (a) $\lsm(\pdist[t]{\gamma}, \pdist[t]{\gamma^*})$ as a function of $t$ and $\gamma$. 
    (b) $\ldsm(\pdist{\gamma}, \pdist{\gamma^*})$ vs. $\gamma-\gamma^*$ for multiple $\gamma^*$, showing non-negligable sensitivity throughout. (c) Mixture weight estimates $\hat\gamma$ from generated data for $\gamma^*=0.6$ under modified linear and cosine noise schedules (see \Cref{sec: Reparam}). The bias and variance of $\hat\gamma$ increase as $L(\gamma^*)$ decreases. (d) $\lsm$ and $\ldsm$ as a function of separation for $\gamma^*=0.8$ and $\gamma=0.9$. $\ldsm$ remains nonnegligible while $\lsm$ rapidly decays.
    }
    \label{fig:gmm prac sens}
\end{figure}


\section{Diffusion sensitivity to GMM mixture weights}

We show that for a bimodal isotropic Gaussian mixture in any dimension, the diffusion score sensitivity index (DSSI) $L(\gamma^*)$ 
is non-negligible. As a result, a diffusion model with a well-trained score generates samples with the correct weights even for extreme values of $\gamma$.  %Our bound is dimension-free and applies to bimodal isotropic mixtures of the following form: $\pdist{\gamma}(\vx) =(1-\gamma) \; \CalN(\vx; \mu_0, \vect I)  + \gamma \;\CalN(\vx; \mu_1, \vect I)$.  


\begin{theorem}[Diffusion sensitivity in bimodal Gaussian mixtures]\label{thm: GMM Sensitivity Emerges}
Consider the family $\pdist{\gamma}(\vx) =(1-\gamma) \; \CalN(\vx; \mu_0, \vect I)  + \gamma \;\CalN(\vx; \mu_1, \vect I)$ and a linear noise schedule $\beta_t=\beta_1 t$ for $\beta_1>0$. 
Let $\alpha_t=\exp\fitparenth{\int_0^t\beta_s \mathrm{d}s}$. % We can remove this if it's unnecessary to redefine this
Then for any $a \in (0,1/2)$
\begin{align}
   \min_{\hat \gamma \in [a,1-a]}\frac{\ldsm(\pdist{\hat\gamma},\pdist{\gamma^*})}{(\hat\gamma-\gamma^*)^2}>\frac{(1-\Phi(\norm{\mu_1-\mu_0}\sqrt{\alpha_1}))\fitparenth{\sqrt{k^2+16}-k}}{7\beta_1\gamma^*(1-\gamma^*)}\exp\fitparenth{-(k^2+k)/4} 
   \label{eq:main_thm}
\end{align}
% \[\min_{\hat \gamma \in [a,1-a]}\frac{\ldsm(\hat\gamma,\gamma^*)}{(\hat\gamma-\gamma^*)^2}>\frac{(1-\Phi(\norm{\mu_1-\mu_0}\sqrt{\alpha_1}))\fitparenth{\sqrt{k^2+16}-k}}{4\beta_1\gamma^*(1-\gamma^*)}\exp\fitparenth{-\frac{1}{4}\fitparenth{k^2+2}}\]
% \becca{
% \begin{align}
% \min_{\hat \gamma \in [a,1-a]}\frac{\ldsm(\hat\gamma,\gamma^*)}{(\hat\gamma-\gamma^*)^2}>&\frac{(1-\Phi(\norm{\mu_1-\mu_0}\sqrt{\alpha_1}))\fitparenth{\sqrt{k^2+16}-k}}{4\beta_1\gamma^*(1-\gamma^*)}\exp\fitparenth{-\frac{1}{4}\fitparenth{k^2+2}}\\
% =&
% \frac{(1-\Phi(\norm{\mu_1-\mu_0}\sqrt{\alpha_1}))\fitparenth{\sqrt{k^2+16}-k}}{4e^{1/2}\beta_1\gamma^*(1-\gamma^*)}\exp\fitparenth{-k^2/4}\\
% >&
% \frac{(1-\Phi(\norm{\mu_1-\mu_0}\sqrt{\alpha_1}))\fitparenth{\sqrt{k^2+16}-k}}{7\beta_1\gamma^*(1-\gamma^*)}\exp\fitparenth{-k^2/4}
% \end{align}
% }
with $k:=\left|\log \frac{1-a}{a}\right|+\left|\log\frac{1-\gamma^*}{\gamma^*}\right|$ and $\Phi(\cdot)$ the CDF of $\mathcal{N}(0,1)$. 

% Let $p_{\mathrm{data}}:=p_{\gamma^*}$ be the data distribution for some $\gamma^*\in [0,1]$. Let $\vs_\btheta(\vx,t)$ be a learned score network with score error \eqref{eq:diffScoreMatchingLoss} $<\delta$ for $\lambda_t=1$.
% \btodo{Our main theorem absolutely 100\% cannot reference a score error defined in the appendix. Also, why sometimes say score error and sometimes say "diffusion score matching loss. essential to be consistent}

% Suppose that $\vs_\btheta(\cdot, t)\equiv\nabla \log \pdist[t]{\hat \gamma}(\cdot )$ for all $t$, for some fixed $\hat \gamma\in [0,1]$. Then, there exists $C=C(\norm{\mu_1-\mu_0}, \gamma^*, \set{\alpha_t})$ such that 
%     \begin{align*}
%         (\hat \gamma-\gamma^*)^2&\leq \frac{\delta}{C}
%     \end{align*}
%     with the minimum value of $C$ achieved at $\gamma^*=0.5$. Equivalently, $L(\gamma^*)\geq C$.

\end{theorem}
% \becca{\cref{thm: GMM Sensitivity Emerges} tells us that when we train a diffusion model to have error less than some value $\delta$, then the corresponding squared error in the mixture weights  }

\emph{Proof Sketch:}
The proof proceeds in three steps. We first decompose $\ldsm(\pdist{\hat\gamma}, \pdist{\gamma^*})$ by mixture component \eqref{eq: mode decomposition} to reduce the problem to lower-bounding the expectations over paths from $\pdist{0}$ and $\pdist{1}$ independently. Second, using the Probability Flow ODE (which is equivalent in probability to the forward process) \citep{song2021scorebasedgenerativemodelingstochastic}, we find an analytical form for $\norm{\score[t]{\hat\gamma}-\score[t]{\gamma^*}}^2$ along paths $\vect{z}_t$ starting from $\pdist{0}$ (\Cref{prop:exact norm expr}). Finally, we identify the interval of $t$ values where this expression is bounded below (\Cref{prop: lower bound over t}), integrate over this region of $t$ (\Cref{prop: Lower bound integral}), and finally integrate over initial conditions $\vect {y} _ {0} $. The bound is dimension-independent because for a given initial condition $\vect {y}_{0} $ starting from e.g. mode $\pdist{0}$, the bound only depends on $\inprod{\vect y_0-\mu_0, \mu_1-\mu_0}/\norm{\mu_1-\mu_0}\sim \mathcal{N}(0,1)$ in any dimension. For full details, see \Cref{Proof of Theorem 1}.
% First, note $\ldsm(\pdist{\hat\gamma}, \pdist{\gamma^*}) = (1-\gamma) \expect[\vect y_t \sim \pdist[t]{0}, t\in (0,1) ]{\|\score[t]{\hat\gamma}(\vect y_t) - \score[t]{\gamma^{*}}(\vect y_t)\|^2} + \gamma \expect[\vect y_t \sim \pdist[t]{1}, t\in (0,1)]{\|\score[t]{\hat\gamma}(\vect y_t) - \score[t]{\gamma^{*}}(\vect y_t)\|^2}$.
% %, where $\pdist[t]{0}$ and $\pdist[t]{1}$ are the marginals of the forward process starting from the unimodal mixture components $p_0$ and $p_1$ respectively \becca{didn't we define these earlier?}. 
% Because $\pdist[t]{0}$ and $\pdist[t]{1}$ are isotropic Gaussians, $\|\score[t]{\hat \gamma}(\vect y_t) - \score[t]{\gamma^{*}}(\vect y_t)\|^2$ has an analytical form conditioned on its initial condition, allowing us to perform a pathwise analysis of each summand.  For a linear noise schedule, we identify the range of values for $t$ where $\|\score[t]{\hat\gamma}(\vect y_t) - \score[t]{\gamma^{*}}(\vect y_t)\|^2$ is large and use this to lower bound $\int_0^1 \|\score[t]{\hat\gamma}(\vect y_t) - \score[t]{\gamma^{*}}(\vect y_t)\|^2\mathrm{d}t$. For full details, see  \Cref{Proof of Theorem 1}.
\qed

\begin{remark}
    Recall that if $\ldsm(\pdist{\gamma},\pdist{\gamma^*})<\delta$ then $(\gamma-\gamma^*)^2 \le \delta/L(\gamma^*)$ by the definition of $L(\gamma^*)$ in \eqref{eq: DSSI}.
    \Cref{thm: GMM Sensitivity Emerges} shows that $L(\gamma^*)$ cannot be too small under mild conditions for two-component Gaussian mixture models. 
    Further, recall that the neural network loss $\lnn$ \eqref{eqn: neural network loss} is an upper bound for $\ldsm$, and so if we train a neural network to an error $\lnn(\pdist{\btheta},\pdist{\gamma})<\delta$ and our neural network learns a score corresponding to the mixture distributions $\{\pdist[t]{\gamma'}\}_{t\in (0,1)}$, then $\ldsm(\pdist{\gamma'},\pdist{\gamma^*})<\delta$. 
    % Together, these statements imply that a well-trained diffusion model will accurately reflect the true mixture weight in the generated distribution. 
    Our theory more directly implies that any score model which closely approximates the true score $\score{\gamma^*}$ in $L^2$ will have small $\ldsm$, and thus its sample distribution will accurately reflect $\gamma^*$.
\end{remark}

\begin{remark}
    If, when generating new samples at inference time, we alter the noise schedule $\{\alpha_t\}$, this may alter the mixture proportions in the generated samples. Specifically, altering the noise schedule changes $\ldsm$ and $L(\gamma^*)$, as noted below \eqref{eq:diffScoreMatchingLoss}.
    % so some choices of noise schedule have $\mathbb{E}_{\vx\sim p^{(\gamma^*)}_t}\|\vs_t^{(\gamma^*)}(\vx)-\vs_t^{(\gamma)}(\vx)\|^2$ small for a large range of values of $\gamma$ and $t$, meaning $L(\gamma^*)$ is smaller. 
    Hence, similar loss under $\ldsm$ for two different noise schedules could potentially correspond to very different guarantees for the accuracy of the recovered mixture weight $\hat\gamma$. Discretizing the reverse process SDE accumulates time-discretization errors regardless of score accuracy, meaning noise schedules that lower $L(\gamma^*)$ may worsen mixture-weight errors in the generative distribution, even for well-trained diffusion models.
    \Cref{sec: GMM Sens Exp} showcases this effect for GMMs. In \Cref{sec:sensitivity-real-data}, we conduct further experiments on real-world data. 
\end{remark}



\paragraph{Effect of parameters on DSSI.} The DSSI $L(\gamma^*)$ depends on the mode separation $\norm{\mu_1-\mu_0}$, true mixture weight $\gamma^*$, and choice of noise schedule $\set{\alpha_t}_t$. First, consistent with the fact that $\ldsm(\pdist{\gamma}, \pdist{\gamma^*})$ is an increasing function of the separation (\Cref{fig: different loss functions}), \Cref{fig: score sensitivity over separation} shows $L(\gamma^*)$ is an increasing function of the mode separation. Second, although \Cref{thm: GMM Sensitivity Emerges} is only stated for linear noise schedules, different continuous noise schedules are equivalent up to time-reparameterization, guaranteeing $\lsm(\pdist[t]{\hat\gamma}, \pdist[t]{\gamma^*})$ is nonnegligible for some range of $t$ under any noise schedule (\Cref{fig: gmm cos heatmap}). Note however that $L(\gamma^*)$ is defined relative to $\ldsm$, which averages $\lsm(\pdist[t]{\hat\gamma},\pdist[t]{\gamma^*})$ over $t\in[0,1]$, meaning a noise schedule which reduces this range of $t$ to a small interval $[a,a+\epsilon]$ scales down $L(\gamma^*)$ roughly proportional to $\epsilon$, regardless of the peak value of $\lsm(\pdist[t]{\hat\gamma},\pdist[t]{\gamma^*})$. 
% This is the mechanism underlying our ablation experiments in \Cref{sec:sensitivity-real-data} and \Cref{sec: GMM Sens Exp}. 
Third, \Cref{fig: score sensitivity ratio over gamma^* and gamma hat} shows that for mode separation $\norm{\mu_1-\mu_0}=9$, $L(\gamma^*)$ is bounded below by 0.37 across all $\gamma^*$, implying that the squared prediction error in $\gamma$ is at most $\approx 2.7$ times larger than the squared diffusion score error. Notably, this bound is dimension-independent, and hence this holds for bimodal isotropic Gaussian mixtures in any dimension. In \Cref{sec: Other Parameters}, we showcase additional experiments on parameters of Gaussian mixtures other than the mixture weight. 


% \textbf{Effect of Separation:} 
% Recall that for fixed $\gamma,\gamma'$, we saw in \Cref{fig: different loss functions} that $\ldsm(\pdist{\gamma}, \pdist{\gamma'})$ increases as a function of separation $\norm{\mu_1-\mu_0}$ (more generally, the Mahalanobis distance between the modes, which reduces to $\norm{\mu_1-\mu_0}$ in the case of isotropic Gaussian modes). \Cref{fig: score sensitivity over separation} shows that similarly, the DSSI is an increasing function of separation for all $\gamma^*$.


% \textbf{Effect of Noise Schedule $\set{\alpha_t}$:} 
% Our theorem guarantees a lower bound on $L(\gamma^*)$ only for linear noise schedules. However, different choices of continuous noise schedule are equivalent up to a rescaling of time; hence, the paths taken along the forward process are equivalent. Therefore, our analysis guarantees that 
% $\lsm(\pdist[t]{\hat\gamma}, \pdist[t]{\gamma^*})$ is nonnegligable for some range of $t$ for any noise schedule. To demonstrate this, \Cref{fig: gmm cos heatmap} shows $\ell_\mathrm{SM}(\pdist[t]{\gamma},\pdist[t]{\gamma^*})$ as a function of $\gamma$ and $t$, as in \Cref{fig: gmm lin heatmap}, but for a squared-cosine noise schedule \citep{nichol2021improveddenoisingdiffusionprobabilistic}.
% Note however that $L(\gamma^*)$ is defined relative to $\ldsm$, which averages $\lsm(\pdist[t]{\hat\gamma},\pdist[t]{\gamma^*})$ over $t\in[0,1]$. Therefore, a noise schedule which reduces this range of $t$ to a small interval $[a,a+\epsilon]$ scales down $L(\gamma^*)$ roughly proportional to $\epsilon$, even if the peak value of $\lsm(\pdist[t]{\hat\gamma},\pdist[t]{\gamma^*})$ over $t$ is large. This is the mechanism underlying our ablation experiments in \Cref{sec:sensitivity-real-data} and \Cref{sec: GMM Sens Exp}.

% \textbf{Effect of $\gamma^*$:} \Cref{fig: score sensitivity ratio over gamma^* and gamma hat} plots $\frac{\ldsm(\pdist{\hat\gamma},\pdist{\gamma^*})}{(\hat\gamma - \gamma^*)^2}$ as a function of $\hat \gamma$ for multiple choices of $\gamma^*$ with separation $\norm{\mu_1-\mu_0}=9$. 
% For any value of $\gamma^*$, we see that empirically $L(\gamma^*)$ is greater than $0.37$, suggesting that the squared prediction error in $\gamma$ is at most $\approx 2.7$ times larger than the squared diffusion score error. We highlight that our analysis is dimension-independent, and hence the sensitivity index is bounded below by $0.37$ 
% for bimodal isotropic Gaussian mixture models in any dimension when $\norm{\mu_1-\mu_0}=9$.


\begin{figure}
\centering
    \begin{subfigure}[htbp]{0.3\linewidth}
    \includegraphics[width=\columnwidth]{images/figure_three_a.pdf}
    \caption{}
    \label{fig: score sensitivity over separation}
    \end{subfigure}
    \hfill    \begin{subfigure}[htbp]{0.3\linewidth}
        \includegraphics[width=\columnwidth]{images/figure_three_b.pdf}
        \caption{}
        \label{fig: gmm cos heatmap}
    \end{subfigure}
    \hfill    \begin{subfigure}[htbp]{0.3\linewidth}
    \includegraphics[width=\columnwidth]{images/figure_three_c.pdf}
        \caption{}
        \label{fig: score sensitivity ratio over gamma^* and gamma hat}
    \end{subfigure}
\caption{
Effect of Gaussian mixture model parameters on score sensitivity in dimension $d=10$.  (a) $L(\gamma^*)$ vs. mode separation for multiple $\gamma^*\leq 0.5$ (by symmetry $L(\gamma^*)=L(1-\gamma^*)$). $L(\gamma^*)$ is an increasing function of separation.
(b) $\lsm(\pdist[t]{\gamma}, \pdist[t]{\gamma^*})$ as a function of $t$ and $\gamma$ for a cosine noise schedule. Compared to a linear noise schedule, the noise schedule only affects the range of $t$ in the forward process at which the score is sensitive to $\gamma$. 
(c) $\ldsm(\pdist{\gamma}, \pdist{\gamma^*})/(\gamma -\gamma^*)^2$ vs. $\gamma$; $L(\gamma^*)\geq 0.27$ for all $\gamma^*$, with the minimum at $\gamma^*=0.5$ . 
}\label{fig:sensitivity_index_gmm}
\end{figure}

\subsection{Related Work} 


\paragraph{Multimodal capabilities of diffusion models.}
The ability of diffusion models to sample multimodal targets varies substantially across common choices of neural architecture \citep{hakemi2025deeper}, training recipes \citep{karras2022elucidating}, and sampling algorithms \citep{roos2026met, zhang2025collapse}. Existing evaluation metrics \citep{sajjadi2018assessing, kynkaanniemi2019improved, djolonga2020precision, naeem2020reliable, alaa2022faithful} are typically structure-agnostic and may fail to reliably detect mode coverage \citep{raisa2025position, stein2023exposing}. 
In contrast, the diffusion score sensitivity index \eqref{eq: DSSI} provides a rigorous framework to reason about multimodal capability.



\paragraph{Learning GMMs with Diffusion Score Matching.} 
Diffusion score matching (DSM) is \textit{in principle} expressive enough to learn Gaussian mixture: gradient descent recovers component means in the balanced isotropic regime\footnote{The mixture model $q\coloneqq \sum_{i=1}^k \frac{1}{k}\, \mathcal{N}(\boldsymbol{\mu}^{(i)},\, \mathbf{I})$ where each component has equal mixture weight and identity covariance} \citep{shah2023learning}, 
learning algorithms can be tailored for density estimation of GMMs with polynomial sample complexity \citep{gatmiry25b, chen2024learning}. Here, \citet{gatmiry25b} develop noise-sensitivity bounds similar in spirit to ours. Unlike the aforementioned articles, our work focuses on the ability of diffusion models to sample from target GMMs rather than on density estimation. 
When the learned score is $L^2$-accurate, the resulting diffusion sampler is provably close to the target in total variation \citep{chen2023sampling, lee2023convergence}. However, this may be insufficient to recover structural parameters that shape the geometry of distributions in low-probability regions. 
\citet{li2024softmixturedenoisingexpressive} suggest diffusion models may fail to sample multimodal targets due to the lack of expressivity of the Gaussian transition steps in the reverse diffusion process. 
Our work complements \citet{li2024softmixturedenoisingexpressive} by identifying a different mechanism for multimodal sampling, the diffusion score sensitivity index \eqref{eq: DSSI}, and by providing additional insights into the role of noise schedules in alleviating mode amplification.

\paragraph{Benefits of Annealing.} 
Classical (unannealed) score matching \citep{10.5555/1046920.1088696} is unsuitable for generative modeling of multimodal targets: the asymptotic efficiency of the score matching estimator relative to the maximum-likelihood estimator is governed by the log-Sobolev (LS) constant \citep{koehler2022statisticalefficiencyscorematching}, which grows exponentially in mode separation for Gaussian mixtures \citep{chen2021dimensionfreelogsobolevinequalitiesmixture}. \citet{song2019generative} proposed annealing over noise scale to handle multimodality, inspired by \citet{neal_annealed_2001}. 
\citet{biroli2024dynamical} and \citet{raya2023symmetry} observe a phase transition at intermediate noise scales linked to the emergence of multimodal structure in practice. 
%Recent work identify annealing as the mechanism behind the ability of diffusion models to sample multimodal distributions. \citet{biroli2024dynamical} observe a speciation phase transition at intermediate diffusion time linked to the emergence of modal structure during sampling, \citet{raya2023symmetry} view the same event as a spontaneous symmetry-breaking bifurcation. 
Amidst the empirical evidence, \citet{qin2024fit} establishes the first formal statistical benefit of annealing in the context of diffusion. \citet{qin2024fit} adapts techniques from the acceleration of Markov Chain mixing to derive a better score matching loss, with the resulting outcome closely resembling the diffusion score matching loss. While \citet{qin2024fit} establishes sample complexity bounds, our work addresses a complementary question of parameter recovery from generated samples. % for the annealed score matching loss, we investigate the ability of generated samples to accurately reflect critical parameters of interest such as mixture weights.

\section{Sensitivity in real datasets}\label{sec:sensitivity-real-data}

We empirically show that mixture weight information is captured in the score functions of intermediate forward process distributions, even outside the Gaussian mixture setting. We additionally explore how altering the noise schedule affects the DSSI and empirically demonstrate that a small DSSI value quantitatively degrades mixture-weight recovery on non-Gaussian data. This is particularly salient to methods like \citep{songddim2021, lu2022dpmsolver, karras2022elucidating} where noise schedules are selected to increase the speed of inference-time sampling; we show that such methods have the potential to amplify some modes over others, augmenting the recent observation by \citet{zhang2025collapse}. 

To study this, we train a simple autoencoder on the ones and eights from MNIST (\citep{lecun2010mnist}). Our target distribution $\pdist{\gamma^*}$ is then the mixture distribution of latent representations of eights and ones ($\pdist{1}$ is the distribution of ones). We set $\gamma^*=0.4$ and use deep ResNets \citep{He2015DeepRL} with a linear noise schedule as our score matching models.

\begin{figure}[ht]
    \centering
    \begin{subfigure}[h]{0.3\columnwidth}
        \centering
        \includegraphics[width=\columnwidth]{images/figure_four_a.pdf}
    \end{subfigure}%
    \begin{subfigure}[h]{0.3\columnwidth}
        \centering
        \includegraphics[width=\columnwidth]{images/figure_four_b.pdf}
    \end{subfigure}%
    \begin{subfigure}[h]{0.3\columnwidth}
        \centering
        \includegraphics[width=\columnwidth]{images/figure_four_c.pdf}
    \end{subfigure}%
    \\
    \begin{subfigure}[h]{0.3\columnwidth}
        \centering
        \includegraphics[width=\columnwidth]{images/figure_four_d.pdf}
        \caption{}
        \label{fig:mnist40}
    \end{subfigure}%
    \begin{subfigure}[h]{0.3\columnwidth}
        \centering
        \includegraphics[width=\columnwidth]{images/figure_four_e.pdf}
        \caption{}
        \label{fig:mnist37}
    \end{subfigure}%
    \begin{subfigure}[h]{0.3\columnwidth}
        \centering
        \includegraphics[width=\columnwidth]{images/figure_four_f.pdf}
        \caption{}
        \label{fig:mnist28}
    \end{subfigure}%
    \caption{
    Influence of noise schedule on mixture weight recovery in MNIST. (a) Linear noise schedule; (b-c) two synthetic noise schedules designed to reduce $L(\gamma^*)$ (see \Cref{sec: Reparam}). Top row: $\lsm(\pdist[t]{\gamma},\pdist[t]{\gamma^*})$ as a function of $t$ and $\gamma$ estimated using trained diffusion models, with $\gamma^*=0.4$. Bottom row: generated samples using each schedule, with estimated $\hat\gamma$ below. Larger $L(\gamma^*)$ yields more accurate mixture weight recovery.
    }
    \label{fig: MNIST experiment}
\end{figure}



\paragraph{Sensitivity at intermediate $t$.} In \Cref{fig: MNIST experiment}, the top row displays the score sensitivity over time and $\gamma$ using trained diffusion models (each with a different noise schedule) to approximate the true score function (see \Cref{sec: MNIST Exp Details} for more details). The first column uses the same linear noise schedule as in \Cref{fig: gmm lin heatmap}. Note that relative to the Gaussian mixture example, $\lsm(\pdist[t]{\gamma}, \pdist[t]{\gamma^*})$ is nonnegligable over a wider range of $t$ values.  
We numerically estimate that the sensitivity index is $\approx 5\times$ larger than in the Gaussian mixture model example, so our theory guarantees a five-fold smaller error in the predicted mixture weight for the same level of error under $\ldsm$ as in the Gaussian setting.

\paragraph{Sensitivity affects parameter recovery.} 
In \Cref{fig:mnist40}, the bottom row shows generated samples from a diffusion model with $T=100$. We estimate the value of $\gamma^*$ using generated samples and a trained classifier, shown below the generated samples. The generated samples accurately reflect the true value of the mixture weight, with only a $1\%$ relative error in estimating $\gamma^*$. 

To demonstrate the effect of the score sensitivity index $L(\gamma^*)$ on the error in estimating $\gamma^*$, we perform an ablation study on the trained diffusion model by modifying the noise schedule to reduce $L(\gamma^*)$ as in \Cref{sec: GMM Sens Exp} (for details see \Cref{sec: Reparam}). \Cref{fig:mnist37} and \Cref{fig:mnist28} show the score sensitivity and generated samples for the modified noise schedules, along with the corresponding estimates $\hat\gamma$ from the samples. When the score sensitivity index reduces by a factor of $10$ (\Cref{fig:mnist28}), we see that the estimate of $\gamma^*$ has a relative error of $31\%$. 
% In other words, even in a real-life distribution, the sensitivity index is a predictive metric for mixture weight recovery. 
Surprisingly, the samples from the modified noise schedules have comparable fidelity to those from the linear noise schedule (\Cref{fig:mnist40}). This is in line with the theoretical and empirical observations that say that diffusion models under errors and accelerated sampling schemes can still accurately learn the geometry of the support of the target distribution \citep{songddim2021, nichol2021improveddenoisingdiffusionprobabilistic, stanczuk2022your,chandramoorthy2026when}. More details on our MNIST experiments can be found in \Cref{sec: MNIST Exp Details}. 

\section{Conclusion}

We investigate the ability of 
diffusion models to recover the relative mode amplitudes of a multimodal target, despite the insensitivity of the score of the target distribution to mixture weights under high separation. 
Our analysis clarifies the impact of annealing: the score of the noisy target distribution across noise scales provides the necessary information for accurate sampling of multimodal distributions. Our experiments indicate that inference-time sampling algorithms can be tuned for diversity and fidelity, suggesting a mechanism for tailoring
sampling algorithms to data distributions arising in scientific applications. 

 
\paragraph{Limitations.} We prove a quantitative relationship between the sensitivity index and the diffusion score error only in the Gaussian mixture setting. Theorem \ref{thm:sensitivity} provides an information theoretic mechanism that cannot be computationally verified since it involves minimax bounds. Further, evaluation of the diffusion score-sensitivity index requires knowledge of the true parameter of interest $\gamma^*$ and the score function of the target distribution $\score[t]{\gamma^*}$. Thus, the application to benchmark datasets requires surrogate generative models in place of the unknown target score.


\bibliographystyle{unsrtnat} %plainnat
\bibliography{refs}

\newpage 
\begin{appendices}
\crefalias{section}{appendix}
\crefalias{subsection}{appendix}

\section{Notation}

{\renewcommand{\arraystretch}{1.5}
\begin{tabular}{ |p{15mm}| p{77mm}| p{35mm} |}
\hline
 Symbol & Definition & In words
 \\
 \hline
 $\mathcal{P}(p_0, p_1)$ 
 & 
 $\set{\pdist{\gamma}: \pdist{\gamma} = \gamma \; p_1+(1-\gamma) \; p_0,\ \gamma \in (0,1)}$ 
 & 
 Family of bimodal mixtures of $p_0$ and $p_1$
 \\
 \hline
 $p_t$ 
 & 
 $\mathrm{d} \vx_t = -\frac{\beta_t}{2} \vx_t \mathrm{d}t + \sqrt{\beta_t}\, \mathrm{d}W_t \text{ and } \vx_0 \sim p \Longrightarrow \vx_t \sim p_t$
 & 
 Distribution of samples at time $t$ in the forward process starting from $p$
 \\
 \hline
 $\lnn(\vs_\btheta, p)$ 
 & 
 $\expect[{ \vx_0 \sim p,\ t,\ \vx_t \mid \vx_0}]{\lambda_t\, \norm{\vs_\btheta(\vx_t, t) - \nabla \log \pdist[t]{\gamma^{*}}(\vx_t \mid \vx_0)}^2}$ 
 & 
 Diffusion model training objective
 \\
 \hline
 $\lsm(p,q)$ 
 & 
 $\expect[{\vx \sim q}]{\norm{\nabla \log p(\vx) - \nabla \log q(\vx)}^2}$ 
 & 
 Classical score matching loss
 \\
 \hline
 $\ldsm(p,q)$ 
 & 
 $\begin{aligned}[t]
   &\expect[{t \in (0,1),\ \vx_t \sim q_t}]{\lambda_t\, \norm{\nabla \log p_t(\vx_t) - \nabla \log q_t(\vx_t)}^2} \\
   &=\expect[{t \in (0,1)}]{\lambda_t\lsm(p,q)}
 \end{aligned}$ 
 & 
 Diffusion score matching loss
 \\
 \hline
 $L(\gamma^*)$ 
 & 
 $\displaystyle \min_{\gamma'} \frac{\ldsm(\pdist{\gamma'},\pdist{\gamma^*})}{(\gamma^* - \gamma')^2}$
 & 
 Diffusion score sensitivity index
 \\
 \hline
 $\vs_\btheta(\vx, t)$ 
 & 
 $\forall\, \btheta \in \bm{\Theta} = \RR^{\#\,\text{Parameters}},\ \vs_\btheta : \RR^d \times [0,1] \rightarrow \RR^d$ 
 & 
 Neural network with parameters $\btheta$
 \\
 \hline
\end{tabular}
}


\section{Comparison of sampling algorithms for multimodal distributions}


Sampling from multimodal distributions using classical Langevin sampling is known to be challenging \citep{NEURIPS2018_c6ede20e, latuszynski2025mcmc}. For a two-component Gaussian mixture, \Cref{fig: different loss functions} shows that as mode separation increases, the sensitivity of the classical score matching loss to the mixture weight decays rapidly to 0 while the diffusion score matching loss does not. As a consequence, sampling algorithms based solely on the target score are unsuitable for sampling from multimodal distributions. \Cref{fig:Langevin_vs_Diffusion} illustrates this effect, comparing classical Langevin sampling to diffusion sampling. 

%a two-component mixture of well-separated Gaussians using classical Langevin sampling is known to be challenging, where samples may be drawn only from one component; in contrast, diffusion score-based sampling accurately samples the true distribution. 


\begin{figure}[ht]
    \centering
    \begin{subfigure}[htbp]{0.19\linewidth}
        \includegraphics[width=\columnwidth]{images/figure_one_a.pdf}
        \caption{}
        \label{fig: true distribution kde}
    \end{subfigure}
    \begin{subfigure}[htbp]{0.19\linewidth}
        \includegraphics[width=\columnwidth]{images/figure_one_b.pdf}
        \caption{}
        \label{fig: langevin sampling kde}
    \end{subfigure}
    \begin{subfigure}[htbp]{0.19\linewidth}
        \includegraphics[width=\columnwidth]{images/figure_one_c.pdf}
        \caption{}
        \label{fig: analytic diffusion kde}
    \end{subfigure}
    \begin{subfigure}[htbp]{0.19\linewidth}
        \includegraphics[width=\columnwidth]{images/figure_one_d.pdf}
        \caption{}
        \label{fig: learned diffusion kde}
    \end{subfigure}
    \caption{Comparison between different sampling approaches on a high-separation Gaussian mixture model. (a) shows a KDE of the true data distribution. (b)-(c) show, in order, KDEs of samples using (b) Langevin sampling (run to time $T=10^5$), (c) diffusion sampling using the analytic score function, and (d) diffusion sampling using a trained neural network. We can see that, unlike Langevin sampling, diffusion sampling accurately reflects the relative masses of the two modes. }
    \label{fig:Langevin_vs_Diffusion}
\end{figure}


% \paragraph{Choice of weighting.}
% The weighting function $\lambda(\cdot)$ is a modeling choice rather than an intrinsic property of the forward process. \citet{song2021scorebasedgenerativemodelingstochastic} propose, 
% \begin{equation}
%     {\lambda_t} \;\propto\; \frac{1}{\expect[{\vx}_0 \sim p_0,\; {\vx}_t \sim p_t(\cdot \mid {\vx}_0)]{\norm{\nabla_{{\vx}_t} \log p_t({\vx}_t \mid {\vx}_0)}^2}},
%     \label{eq:lambda-spec}
% \end{equation}
% which normalizes each noise level so that the expected squared conditional-score norm contributes evenly to the outer integral in \eqref{eq:dsm-loss}. For the VP-SDE, substituting ${\vx}_t = \sqrt{{\alpha_t}} {\vx}_0 + \sqrt{1-\alpha{t}} \vect{z}$ where $\vect{z} \sim \mathcal{N}(0, \vect{I})$ gives, 
% \[
% \expect{\norm{\nabla \log p_t({\vx}_t \mid {\vx}_0)}^2} = (1-\alpha_t)\cdot t
% \]
% thus, \eqref{eq:lambda-spec} reduces, up to a constant $d$, to 
% \begin{equation}
%     {\lambda_t} = 1-{\alpha_t} \label{eq:lambda-vpsde}
% \end{equation}
% The choice \eqref{eq:lambda-vpsde} downweights errors near $t=0$, since ${\lambda_t} \rightarrow 0$.



% \section{Theoretical background on high-dimensional sampling}
% \subsection{Review: Langevin Sampling and the Log-Sobolev Constant}
% A similar sampling scheme to Score-Based Diffusion is classical Langevin sampling. Assume we have some Gibbs distribution $p(\vx)\propto \exp(E(\vx))$ for which we only have access to the energy function $E$. Then, we can define a continuous-time Markov process 
% \[\mathrm{d} \vx = \frac{1}{2}\nabla_{\vx}E(\vx)\mathrm{d} t+ \mathrm{d}\vect W_t\]

% If we begin with samples from a distribution $q$ and let $q_t$ be the distribution of this process at time $t$, it holds that 
% \[\vect{KL}(q_t| p)\leq e^{-t/C_{LS}(p)} \vect{KL}(q|p)\]
% where $\vect{KL}$ is the Kullback–Leibler divergence and the {\bf Log-Sobolev Constant} $C_{LS}(p)$ is the smallest constant such that for all distributions $\pi$,
% \begin{align*}
%     \vect{KL}(\pi|p)&\leq C_{LS}(p)\mathcal{I}(\pi|p)
%     \\&= C_{LS}(p)\expect[p]{\inprod{\nabla \log \frac{\pi}{p}, \nabla \frac{\pi}{p}}}
% \end{align*}

% where $\mathcal{I}(\pi|p)$ is the Fisher information of $\pi$ given $p$. Thus, Langevin Diffusion is ergodic, with rate determined by the Log-Sobolev constant. References for these bounds can be found in \cite{vanHandelRamon}. Importantly, distributions $p$ for which the Log-Sobolev constant is large require a long burn-in time for the distribution of generated points converge to $p$. 
% %The implications of this will be discussed further in \Cref{GMM Surprise}.

% To prevent confusion, we will refer to the dynamics described in equation \eqref{eqn:rev} as {\bf Annealed Langevin Diffusion} and the dynamics described here as {\bf Classical Langevin Sampling} for the remainder of the paper. Much of the remainder of the manuscript centers around the fact that many distributions which are challenging for Classical Langevin Diffusion are reasonable for Annealed Langeivn Diffusion.

% \subsection{Review: Classical score matching for Langevin Diffusion}\label{Sec: SM Loss}
% Often, in practice, we don't have access to this true energy function $E$, and instead seek to approximate it. This can be done for parameter estimation or for sampling via classical Langevin Diffusion (\cite{10.5555/1046920.1088696}). 

% \begin{definition}
%     For two distributions $\pi$ and $\rho$ on $\RR^d$, we define the {\bf score matching Loss} of $\pi$ with respect to $\rho$ as
%     \begin{align*}
%     D_{\mathrm{SM}}(\pi|\rho):&=\expect[\vx\sim \rho]{\norm{\nabla \log \rho(\vx)- \nabla \log \pi(\vx)}^2}
%     \\
%     &=\expect[\vx\sim \rho]{\mathrm{Tr}\;\nabla^2\log \pi(\vx)+ \norm{  \nabla \log \pi(\vx)}^2} + C_\rho
%     \end{align*}
%     where $C_p$ is independent of $\pi$. This second form of the equation is important as it allows for the calculation of this loss without knowledge of $\nabla \log \rho$, although we will mostly focus on the first form for this paper.
%     Note that this is not symmetric, as the samples in the inner expectation are drawn from $\rho$.
% \end{definition}

% If we want to generate samples from a distribution $p$ we can only approximate from data, then there exist two challenges: 
% \begin{enumerate}
%     \item Estimation of $\nabla \log p$
%     \item Using the estimated score to sample from $p$ via Classical Langevin Diffusion
% \end{enumerate}

% While the latter obstacle is intrinsic to the properties of $p$ as discussed before, the former can also be challenging for certain families of distributions, as we will discuss later.

\section{Minimax upper and lower bounds}
\label{sec:lowerUpperBounds}
Here we provide a more detailed proof of Theorem \ref{thm:sensitivity} and provide additional remarks on the Theorem and the underlying assumptions. We begin with a remark on score-restricted estimators (Definition \ref{def:scoreEstimator}) and explain the motivation for their definition.

In minimax theory, and in particular, Fano's method, the accuracy of an estimator for a parameter learned from a finite number of data samples is inherently limited by how well the data samples can distinguish between two parameter values. This distinguishability is often measured in terms of the KL divergence between the data distributions at two nearby parameter values. At first glance, the data samples given to a diffusion model appear to be the target samples, and as noted previously, the data distributions with well-separated modes do not distinguish between nearby values of weight parameters. Further, note
% the noised samples also do not have sufficiently informative data distributions since
$\mathrm{KL}(\pdist[t]{\gamma}||\pdist[t]{\gamma'}) \leq
\mathrm{KL}(\pdist{\gamma}||\pdist{\gamma'}).$ As a result, it might seem that the generated samples from the denoising process do not accurately recover the parameter (particularly, given the role that the KL divergence plays in Fano's inequality \citep{binyu}), but our empirical results disagree with this expectation. 
In our empirical results, we see that the noising/denoising process affects parameter recovery compared with score approximation followed by Langevin dynamics. A key idea that helps resolve this mismatch is to recognize that the data a diffusion model uses to accurately estimate parameters is the noisy score (approximated by a neural network), not the noisy sample data. In other words, the generated samples, from the reverse dynamics, are a function of the noisy scores (interpolated by neural networks). These neural network interpolations at some time $t$ can have a distribution that distinguishes between parameter values more than at time $0$. In that sense, the time $t$ scores contain sufficient information to generate samples corresponding to a nearby parameter value, while the time $0$ score (target score) may not. 

Our score-restricted estimators define the notion of recovering $\gamma^*$ from ``score data'' rather than sample data. Such a notion models parameter estimation from the generated samples in score-based generative models where the score field is first approximated and then used for sampling. To distinguish Langevin sampling from diffusion models, we define two separate classes of estimators: one class has access to only a target score oracle and the other class has access to a score oracle at an intermediate time along the noising process. Hence, the estimator classes $\Gamma_{\mathrm{score},t}$ and $\Gamma_{\mathrm{score},0}$ differ in the data distributions given to them, which are respectively, $(\qdist[t]{\gamma})$ and $(\qdist[0]{\gamma}).$ In Lemma \ref{remark:dataProcessing}, we show that the oracle model $\qdist[t]{\gamma}$ cannot be reconstructed in a $\gamma$-independent manner from the oracle manner $\qdist[0]{\gamma},$ and this induces a separation between the two estimator classes.

\begin{remark}    \label{remark:score-restricted-estimators}
 While in Theorem \ref{thm:sensitivity}, we are only comparing $\Gamma_{\mathrm{score},0}$ and a set of estimators that use score data at one other time $t,$ the result continues to hold if instead under a specific redefinition of the two estimator classes. Let $t$ still be a time at which the diffusion sensitivity definition holds for the values of $\epsilon$ and $\delta$ defined in Theorem \ref{thm:sensitivity}. The estimator class $\Gamma_{\mathrm{score},0}$ contains any sampling scheme that does not use samples from $\qdist[t']{\gamma}$, for any time except $t' = 0.$ On the other hand, the estimators in $\Gamma_{\mathrm{score},0}$ must include sampling (denoising) steps that access score data from $\qdist[t]{\gamma}.$  


\end{remark}

\begin{remark}
    \label{remark:cameronMartin} To compute the KL divergence, $\mathrm{KL}(\qdist[t]{\gamma}||\qdist[t]{\gamma'})$, we first recognize that both Gaussian measures are on $L^2(\pdist[t]{\gamma}).$ Hence, by Cameron-Martin formula, the Radon-Nikodym derivative, $d\qdist[t]{\gamma}/d\qdist[t]{\gamma'}.$ We can take expectations with respect to $\qdist[t]{\gamma}$ to arrive at the formula $\mathrm{KL}(\qdist[t]{\gamma}||\qdist[t]{\gamma'}) = \|\score[t]{\gamma} - \score[t]{\gamma'}\|^2_{L^2(\pdist[t]{\gamma}}/(2\sigma_t^2).$ This resembles the KL divergence between finite dimensional Gaussian distributions with the same covariance. 
\end{remark}

\begin{lemma}
\label{remark:dataProcessing} There is no $\gamma$-independent post processing of score observations at time 0, $\hat{\mathbf s}^{(\gamma)}_0,$ that can reproduce  $\hat{\mathbf s}^{(\gamma)}_t.$
\end{lemma}
\begin{proof}
We can prove why by contradiction. If there exists a $\gamma$-independent processing of $\qdist{\gamma}$ that produces the distribution $\qdist[t]{\gamma},$ the data processing inequality gives $\mathrm{KL}(\qdist[t]{\gamma}||\qdist[t]{\gamma'}) \leq \mathrm{KL}(\qdist{\gamma}||\qdist{\gamma'}).$ But, we have assumed that at time $t,$ our mixture family is such that $\mathrm{KL}(\qdist{\gamma}||\qdist{\gamma'}) \leq \epsilon_t^2/4 \leq (1/128) (\gamma - \gamma')^2/A_t^2 \leq (1/128)\mathrm{KL}(\qdist[t]{\gamma}||\qdist[t]{\gamma'}).$ Thus, the oracle observations and the statistical experiment that can be performed with the model $\qdist[t]{\gamma}$ cannot be simulated from having access to the oracle, $\qdist{\gamma}.$ 
\end{proof}

Before we give the proof of Theorem \ref{thm:sensitivity}, we recapitulate a corollary of Generalized Fano's method from Lemma 3 of \citep{binyu} and define covering entropy, which is used in Theorem \ref{thm:sensitivity}.
\begin{lemma}{[Lack of identifiability of mixture weights]} Let $G:= \{\gamma_1, \cdots, \gamma_N\} \in (0,1)^N$ be a set of $N\geq 16$ values of mixture weights such that $|\gamma_i - \gamma_j| \geq \varepsilon,\: \: i \neq j, \; 1\leq i,j \leq N.$ Then, for the class of score-restricted estimators (\cref{def:scoreEstimator}), $\Gamma_{\mathrm{score},0},$ with oracle access to the model, $\qdist{\gamma^*}$,   if $\mathrm{KL}(\qdist{\gamma_i}||\qdist{\gamma_j}) \leq \log N/4,$ we have
\begin{align}
\inf_{\hat{\gamma} \in \Gamma_{\mathrm{score},0}} \max_{\gamma^* \in G} \mathbb{E}_{\qdist{\gamma^*}} [|\hat{\gamma} - \gamma^*|^2] \geq \varepsilon^2/8.
\label{eq:fano}
\end{align}
\label{lem:fano}
\end{lemma}

We note that the proof follows from \citep[Lemma 3]{binyu}, which is an application of Fano's inequality. To get the error in the squared norm, $\mathbb{E}[|\hat{\gamma} - \gamma^*|^2,$ as we have in \eqref{eq:fano}, we need to additionally use Markov's inequality: $\mathbb{E}[|\hat{\gamma} - \gamma|^2 \geq (\varepsilon/2)^2 \mathbb{P}(|\hat{\gamma} - \gamma^*| > \varepsilon/2) \geq (\varepsilon/2)^2 \: (1/2).$ The second inequality directly follows from the proof of Lemma 3 of \citep{binyu}. The main difference here is that the data used by the estimator corresponds to the score (\cref{def:scoreEstimator}), chosen to exemplify the diffusion model setting.

\begin{definition}{[$\epsilon$-covering entropy]}\label{def:cover-entropy}
Let $O_t$ be a finite open cover of $[0,1],$ with centers $\gamma_1, \cdots, \gamma_{N_t},$ such that for any $\gamma \in (0,1),$ $\min_j \mathrm{KL}(\qdist[t]{\gamma}||\qdist[t]{\gamma_j}) < \epsilon^2.$ Then, $V_{\epsilon,t} := \log |O_t|$ is called the $\epsilon$-covering entropy. 
    
\end{definition}

The main novelty in the proof of Theorem \ref{thm:sensitivity} is the setup of score oracle-based estimators. Note that the above standard definitions of packing and covering entropies extend verbatim to our infinite-dimensional statistical models.


Secondly, while annealing has been known to be statistically beneficial \citep{koehler2022statisticalefficiencyscorematching, qin2024fit}, our result proves better minimax bounds for parameter recovery, rather than density estimation. 
To complete the upper bound in \eqref{eq:upperBound}, we apply directly Theorem 2 \citep{yangBarron}. To obtain the lower bound, we use Lemma \ref{lem:fano} with a possibly different packing set of size $N(\gamma^*)$ around each $\gamma^*.$ Using the definition of $\epsilon$-sensitivity and assuming that $\log N(\gamma^*) \geq \max\{\log 16, \sigma_t^2\: V(\epsilon_{t}) \}$ at all $\gamma^*,$ we obtain that for any $\gamma, \gamma' \in S(\gamma^*),$ $\mathrm{KL}(\qdist{\gamma}||\qdist{\gamma'}) \leq \delta = \epsilon_{t}^2/(4 \sigma_t^2) \leq (\log N(\gamma^*)/V(\epsilon_{t}))  \epsilon_{t}^2/4 = \log N(\gamma^*)/4,$ since $\epsilon_{t}^2 \: \sigma_t^2 = V(\epsilon_{t}).$ As a result, the conditions of Lemma \ref{lem:fano} apply at every $\gamma^*.$



\begin{comment}
    

\subsection{Connecting notions of sensitivity}\label{sec:connecting sensitivity notions}

Here we connect the DSSI $L(\gamma)$ to the information-theoretic notion of sensitivity used in \cref{thm:sensitivity}. Specifically, we show that under some simplifying assumptions, there exist a constant $\sigma^2$ and a time $t^*$ such that for each $\gamma^*$ and $c>0$, there exists a neighborhood $U$ of $\gamma^*$ such that for all $\gamma\in U$,
\[(\gamma -\gamma^*)^2\leq \frac{2\sigma^2}{L(\gamma^*)-c}\mathrm{KL}(\qdist[t^*]{\gamma}||\qdist[t^*]{\gamma^*}) \]

Note this is a weaker condition than \cref{assumption:continuity}, as it only holds locally around each $\gamma^*$. However, it suggests that a family of distributions with large DSSI $L(\gamma^*)$ for all $\gamma^*$ is likely to satisfy \cref{assumption:continuity} for a reasonably small value of $A_{t}$. \cref{thm:sensitivity} requires that the family of distributions exhibits an $\epsilon$-sensitivity for specified values of $\epsilon$ and $\delta$, which is easier to satisfy when $A_{t^*}$ is smaller. In other words, we expect that for a family of distributions with large DSSI for all values of $\gamma^*$, the information-theoretic obstruction caused by the lack of sensitivity of the score at high separation is removed when using score observations at lower separations.


Fix $m\in \mathbb{N}$ . As in \cref{def:scoreEstimator}, we let $\mathcal{S}_t:=\{\score[t]{\gamma^*}(\vx^{(1)}_t), \cdots, \score[t]{\gamma^*}(\vx^{(m)}_t)\}$ where $\vx^{(i)}_t \sim \pdist[t]{\gamma^*}$ are noised samples at time $t$, let $\qdist[t]{\gamma^*}$ be the distribution of the score data, ($\qdist[t]{\gamma^*} = (\score[t]{\gamma^*})_\sharp \pdist[t]{\gamma^*}$), and use $\mathbb{E}_{\gamma^*,t}$ as a shorthand for expectations with respect to $(\qdist[t]{\gamma^*})^{\otimes m}$ (note that $\mathcal{S}_t\sim (\qdist[t]{\gamma^*})^{\otimes m}$). We make two assumptions to simplify our analysis:
\begin{assumption}\label{assump:t exists}
    There exists a time $t^*$ such that for all $\gamma$ and $\gamma'$,
    \[L(\gamma)  \leq \frac{\lambda_{t^*}\,\lsm(\pdist[t^*]{\gamma'}, \pdist[t^*]{\gamma})}{(\gamma' -\gamma)^2}.\]
\end{assumption}

\begin{remark}
    Even for a distribution with large DSSI, it is still possible that for different values of $\gamma$ and $\gamma'$, $\lambda_t\,\lsm(\pdist[t]{\gamma'}, \pdist[t]{\gamma})$ is large for different values of $t$. So, the existence of a fixed $t^*$ such that \cref{assump:t exists} holds across all $\gamma$ and $\gamma'$ is not guaranteed. Empirically, all distributions we tested satisfied \cref{assump:t exists}. 
    % Andrew: I wanted to include this possible way to prove the existence of the t^*. However, as it exists here, the t^* would depend on gamma^*, and I wasn't sure of a good way to get rid of that constraint
    
    % A sufficient, but not neccesary, condition for \cref{assump:t exists} to hold is $\frac{\lambda_t\; \lsm(\pdist[t]{\gamma'}, \pdist[t]{\gamma})}{(\gamma'-\gamma)^2}$ is (quasi-)convex as a function of $\gamma'$ for all fixed $t$ and (quasi-)concave as a function of $t$ for all fixed $\gamma'$. 
    % One can then apply Sion's Minimax Theorem \cite{pjm/1103040253} to the definition of the DSSI, using
    % \[L(\gamma^*) =\inf_\gamma \frac{\int_0^1 \lambda_t\, \lsm(\pdist[t]{\gamma}, \pdist[t]{\gamma^*})\mathrm{d}t}{(\gamma-\gamma^*)^2}\leq \inf_\gamma\max_t\frac{ \lambda_t\, \lsm(\pdist[t]{\gamma}, \pdist[t]{\gamma^*})}{(\gamma-\gamma^*)^2}\]

    % Andrew: More detailed proof (slightly old, would need to be cleaned up)
    % By the definition of $L(\gamma^*)$, we have 
    % \[L(\gamma^*) =\inf_\gamma \frac{\int_0^1 \lambda_t\; \lsm(\pdist[t]{\gamma}, \pdist[t]{\gamma^*})\mathrm{d}t}{(\gamma-\gamma^*)^2}\leq \inf_\gamma\max_t\frac{ \lambda_t\; \lsm(\pdist[t]{\gamma}, \pdist[t]{\gamma^*})}{(\gamma-\gamma^*)^2}\]
    
    % Then, the assumption on $\frac{\lambda_t\; \lsm(\pdist[t]{\gamma}, \pdist[t]{\gamma^*})}{(\gamma-\gamma^*)^2}$ allows us to apply Sion's Minimax Theorem \cite{pjm/1103040253} to get that 
    % \[\inf_\gamma\max_t\frac{ \lambda_t\; \lsm(\pdist[t]{\gamma}, \pdist[t]{\gamma^*})}{(\gamma-\gamma^*)^2}=\max_t\inf_\gamma\frac{ \lambda_t\; \lsm(\pdist[t]{\gamma}, \pdist[t]{\gamma^*})}{(\gamma-\gamma^*)^2}.\]
    % Letting $t^*$ be the point where the outer maximum is achieved, we have that for all $\gamma$
    % \[L(\gamma^*)(\gamma-\gamma^*)^2 \leq   \lambda_{t^*}\; \lsm(\pdist[t^*]{\gamma}, \pdist[t^*]{\gamma^*})\]
\end{remark}

\begin{assumption}\label{assump:est exists}
    For the $t^*$ from \cref{assump:t exists}, there exists an unbiased estimator $\hat\gamma$ using the score information $\CalS_{t^*}$ such that, if $\gamma^*$ is the true value of the mixture weight $\gamma$,
    \[\mathbb{E}_{\gamma^*,t^*} \fitbracket{\lambda_t\,\lsm\fitparenth{\pdist[t^*]{\hat\gamma(\CalS_{t^*})}, \pdist[t^*]{\gamma^*}}}<\sigma^2\]
    for some $\sigma>0$. 
\end{assumption}

\begin{remark}
    Note that because $\gamma$ lies in a compact set, \cref{assump:est exists} holds trivially for some $\sigma$ as long as there exists any unbiased estimator of $\gamma^*$ using the data $\CalS_{t^*}$ and $\qdist[t^*]{\gamma}$ has finite variance for all $\gamma$. If this estimator is asymptotically efficient, then we additionally have that $\sigma^2$ is bounded as a function of $m$ as $m$ grows large (note that $ \mathbb{E}_{\gamma^*,t}$ hides a dependence on $m$).
\end{remark}



Fix the true value $\gamma^*$.
If we let $f_t(\gamma):= \lambda_t\, \lsm(\pdist[t]{\gamma}, \pdist[t]{\gamma^*})$, direct calculation shows 
\[f_t(\gamma^*)=f'_t(\gamma^*)=0,\quad\quad f_t''(\gamma^*)=\expect[\vect x_t \sim \pdist[t]{\gamma^*}]{\norm{\partial_\gamma\score[t]{\gamma}(\vx_t)|_{\gamma=\gamma^*}}^2}.\]

Therefore,
\[ \lambda_t\, \lsm(\pdist[t]{\gamma}, \pdist[t]{\gamma^*})= \frac{1}{2}\expect[\vect x_t \sim \pdist[t]{\gamma^*}]{\norm{\partial_\gamma\score[t]{\gamma}(\vx_t)|_{\gamma=\gamma^*}}^2}(\gamma-\gamma^*)^2 + o(|\gamma-\gamma^*|^2).\]

Using the $t^*$ from \cref{assump:t exists}, we have 
\[L(\gamma^*)\leq  \frac{1}{2}\expect[\vect x_{t^*} \sim \pdist[{t^*}]{\gamma^*}]{\norm{\partial_\gamma\score[{t^*}]{\gamma}(\vx_{t^*})|_{\gamma=\gamma^*}}^2} + o(1).\]

Next, for the estimator $\hat\gamma$ from \cref{assump:est exists}, using \cref{assump:t exists} we have
\[\sigma^2 > \mathbb{E}_{\gamma^*,t^*} \fitbracket{ \lambda_t\, \lsm(\pdist[t]{\hat{\gamma}(\CalS_{t^*})}, \pdist[t]{\gamma^*})} \geq L(\gamma^*)\mathbb{E}_{\gamma^*,t^*} \fitbracket{ (\gamma -\gamma^*)^2}.\]
Note that the variance of $\hat\gamma$ is given by $\mathrm{Var}(\hat{\gamma})=\mathbb{E}_{\gamma^*,t^*}[(\hat{\gamma}(\CalS_{t^*})-\gamma^*)^2]$, so rearranging gives 
\[\mathrm{Var}(\hat{\gamma}) < \frac{\sigma^2}{L(\gamma^*)} \]

If we consider the Fisher information for $(\qdist[t]{\gamma^*})^{\otimes m}$
\[I_{t^*}^{(m)}(\gamma^*) = -\;\mathbb{E}_{\gamma^*,t^*}\fitbracket{ \partial_\gamma^2\log  (\qdist[t^*]{\gamma^*})^{\otimes m}}, \]
then the Cramer-Rao Inequality gives us that 
\[I_{t^*}^{(m)}(\gamma^*)\geq \frac{1}{\mathrm{Var}(\hat{\gamma})} > \frac{L(\gamma^*)}{\sigma^2}.\]

Finally, because
\[\mathrm{KL}(\qdist[t^*]{\gamma}||\qdist[t^*]{\gamma'})= \frac{I_{t^*}^{(m)}(\gamma^*)}{2}(\gamma -\gamma^*)^2 + o(|\gamma-\gamma^*|^2)\]
we have
\[\mathrm{KL}(\qdist[t^*]{\gamma}||\qdist[t^*]{\gamma'}) \geq \frac{L(\gamma^*)}{2\sigma^2}(\gamma -\gamma^*)^2 + o((\gamma-\gamma^*)^2)=\frac{L(\gamma^*)+o(1)}{2\sigma^2}(\gamma -\gamma^*)^2\]

Therefore, for all $c>0$, there exists a neighborhood $U$ of $\gamma^*$ such that for all $\gamma\in U$,
\[\mathrm{KL}(\qdist[t^*]{\gamma}||\qdist[t^*]{\gamma'}) \geq \frac{L(\gamma^*)-c}{2\sigma^2}(\gamma -\gamma^*)^2 \]
\qed

%  We can relate the information-theoretic notion of sensitivity in Assumption \ref{definition:diffusionSensitivity} and the sensitivity of the function $\score[t]{\gamma}$. The main work is to show that an increased parameter sensitivity, $\partial_\gamma\mathbb{E}_{\gamma, t}\score[t]{\gamma},$ when compared to its value at $t=0,$ results in a corresponding increase in the Fisher information at time $t$, $I_t(\gamma):= -\mathbb{E}[\partial^2_\gamma\log \qdist[t]{\gamma}].$ To do this, we assume that the score data variance is uniformly bounded by $\sigma^2$. Then, using Cramer-Rao inequality, we obtain $I_t(\gamma) \geq \mathbb{E}[\partial_\gamma\score[t]{\gamma}]/\sigma^2$.  Since 
% $\mathrm{KL}(\qdist{\gamma}||\qdist{\gamma'}) = (\gamma - \gamma')^2\: I(\gamma)/2 + o((\gamma - \gamma')^2),$ and we know that $I_t(\gamma)$ increases with $t$, the upper bound on the KL increases and therefore the minimax lower bound above decreases. In other words, the information-theoretic obstruction caused by the lack of sensitivity of the score at high separation is removed when using score observations at lower separations.   
\end{comment}

\section{Diffusion process on Gaussian mixture models}
\subsection{Mixture Family Preliminaries}
%\label{subsec:mixture-family-def}

Let $\pdist{0}$ and $\pdist{1}$ be two fixed probability distributions on $\mathbb{R}^d$. We recall the family of mixed probability densities
\begin{equation*}
    \mathcal{P}(\pdist{0},\pdist{1}) \;\coloneqq\; \big\{\, \pdist{\gamma} \coloneqq (1-\gamma)\,\pdist{0}({\vx}) + \gamma\,\pdist{1}({\vx}) \;\big|\; \gamma \in [0,1] \,\big\}.
\end{equation*}
The set of noisy marginals  at noise level $t$ produced by the forward process \eqref{eq:forward-sde} applied to samples ${\vx}_0 \sim \pdist{\gamma}$ is denoted $\{\pdist[t]{\gamma}\}_{t\in [0,1]}$. 
%$$\mathcal{P}(\pdist[t]{0},\pdist[t]{1}) = \{\pdist[t]{\gamma} := (1-\gamma)\pdist[t]{0} + \gamma \pdist[t]{1}\}$ for $t\in[0,1]$. 
The forward process preserves the mixture structure: the noisy marginal $\pdist[t]{\gamma}$ is itself a mixture of the noisy component marginals $\pdist[t]{0}$ and $\pdist[t]{1}$, with the same mixture weight $\gamma$. The family of noisy mixtures and the noisy family of mixtures therefore coincide.
%\subsubsection{}\label{proof: mixture-closure}
\begin{lemma}[Mixture closure under forward VP-SDE]
\label{lemma:mixture-closure}
For every $\gamma \in (0,1)$ and every $t \in [0,1]$, 
\begin{equation}
    \pdist[t]{\gamma} ({\vx}) = (1-\gamma) \pdist[t]{0}({\vx}) + \gamma \pdist[t]{1}({\vx}) 
    \label{eq:mixture-closure}
\end{equation}
\end{lemma}
\begin{proof}
Let $k_t(\vect{y} \mid {\vx}) := \CalN(\vect{y};\, \sqrt{\alpha_t}\,{\vx},\, (1-\alpha_t)\,\vect{I})$ denote the forward conditional density~\eqref{eq:forward-kernel}. Since $k_t$ does not depend on the initial distribution, the marginal at noise level $t$ is obtained by a linear convolution against the initial density:
\[
    \pdist[t]{\gamma}({\vx}) \;=\; \int_{\RR^d} k_t({\vx} \mid {\vx}_0)\,\pdist{\gamma}({\vx}_0)\,\mathrm{d}{\vx}_0.
\]
Substituting the definition of the mixture and applying the linearity of integration,
\begin{align*}
    \pdist[t]{\gamma}({\vx})
    \;&=\; (1-\gamma) \int_{\RR^d} k_t({\vx} \mid {\vx}_0)\,\pdist{0}({\vx}_0)\,\mathrm{d}{\vx}_0
    \;+\; \gamma \int_{\RR^d} k_t({\vx} \mid {\vx}_0)\,\pdist{1}({\vx}_0)\,\mathrm{d}{\vx}_0
    \\
    \;&=\; (1-\gamma)\,\pdist[t]{0}({\vx}) + \gamma\,\pdist[t]{1}({\vx}).
\end{align*}
\end{proof}

A useful consequence of \Cref{lemma:mixture-closure} is that samples from $\pdist[t]{\gamma}$ can be drawn by first sampling a component label $c \sim \mathrm{Ber}(\gamma)$ and then drawing $\vect{y}_t \sim \pdist[t]{c}$, a component-wise decomposition that simplifies subsequent analysis.

\subsubsection{The score of the noisy mixture}

Differentiating $\log \pdist[t]{\gamma}$ and applying \Cref{lemma:mixture-closure} gives the score of the noisy mixture as a pointwise convex combination of the component scores,
\begin{equation}
    \nabla_{\vx} \log \pdist[t]{\gamma}({\vx})
    \;=\; w_0^t({\vx})\,\nabla_{\vx} \log \pdist[t]{0}({\vx}) \;+\; w_1^t({\vx})\,\nabla_{\vx} \log \pdist[t]{1}({\vx}),
    \label{eq:mixture-score-app}
\end{equation}
where the weights are the EM responsibilities --- the posterior probabilities under the generative model that ${\vx}$ was produced by component $0$ or component $1$, respectively \cite[\S 9.2]{10.5555/1162264},
\begin{equation}
    w_0^t({\vx}) \;\coloneqq\; \frac{(1-\gamma)\,\pdist[t]{0}({\vx})}{\pdist[t]{\gamma}({\vx})}, \qquad
    w_1^t({\vx}) \;\coloneqq\; \frac{\gamma\,\pdist[t]{1}({\vx})}{\pdist[t]{\gamma}({\vx})},
    \label{eq:mixture-weights-app}
\end{equation}
and satisfy $w_0^t({\vx}) + w_1^t({\vx}) = 1$.

\subsubsection{Specialization: Gaussian components}

When the components are Gaussian, $\pdist{c} = \mathcal{N}(\boldsymbol{\mu}_c^*, \boldsymbol{\Sigma}_c^*)$ for $c \in \{0,1\}$, the noisy marginals remain Gaussian, $\pdist[t]{c} = \mathcal{N}(\boldsymbol{\mu}_{c,t}^*, \boldsymbol{\Sigma}_{c,t}^*)$, and the score \eqref{eq:mixture-score-app} takes the explicit form
\begin{equation}
    \nabla_{\vx} \log \pdist[t]{\gamma}({\vx}) \;=\; \sum_{c \in \{0,1\}} w_c^t({\vx})\, (\boldsymbol{\Sigma}_{c,t}^*)^{-1} (\boldsymbol{\mu}_{c,t}^* - {\vx}),
    \label{eq:gaussian-mixture-score}
\end{equation}
a responsibility-weighted sum of precision-scaled displacements $(\boldsymbol{\Sigma}_{c,t}^*)^{-1}(\boldsymbol{\mu}_{c,t}^* - {\vx})$. The responsibility weights \eqref{eq:mixture-weights-app} admit the explicit closed form
\begin{equation}
    w_c^t({\vx}) \;=\; \frac{\pi_c\, \mathcal{N}\!\left({\vx};\, \boldsymbol{\mu}_{c,t}^*,\, \boldsymbol{\Sigma}_{c,t}^*\right)}{\sum_{c' \in \{0,1\}} \pi_{c'}\, \mathcal{N}\!\left({\vx};\, \boldsymbol{\mu}_{c',t}^*,\, \boldsymbol{\Sigma}_{c',t}^*\right)},
    \label{eq:responsibility-weights}
\end{equation}
with mixture weights $\pi_0 = 1-\gamma$ and $\pi_1 = \gamma$. The numerator and denominator together carry information about $\gamma$, the Mahalanobis distance from ${\vx}$ to each component mean, and the Gaussian normalizing constants $\det(\boldsymbol{\Sigma}_{c,t}^*)^{-1/2}$.

\subsubsection{DDPM-style forward process} 

\begin{lemma}\label{lemma:mix_over_time}
For $\pdist{0}:=\mathcal{N}(\vx;\mu_0,\vect I)$ and $\pdist{1}:=\mathcal{N}(\vx;\mu_1,\vect I)$ denote two individual component distributions. For a variance-preserving forward process with noise schedule $(\bar{\alpha}_t)_{t \in [0,T]}$, the parameters of the mixture modes at time $t$ are
\begin{equation}
    \boldsymbol{\mu}_{c,t}^* \;=\; \sqrt{\bar{\alpha}_t}\, \boldsymbol{\mu}_c^*, \qquad
    \boldsymbol{\Sigma}_{c,t}^* \;=\; \bar{\alpha}_t\, \boldsymbol{\Sigma}_c^* + (1 - \bar{\alpha}_t)\, \mathbf{I},
    \label{eq:ddpm-perturbation}
\end{equation}
Hence, the noisy marginal distributions are,     
    \begin{align*}
        \pdist[t]{0}(\vx) = \mathcal{N}(\vx;\sqrt{\alpha_t} \mu_0, \vect I),
          \qquad 
        %\\
        \pdist[t]{1}(\vx) = \mathcal{N}(\vx;\sqrt{\alpha_t} \mu_1, \vect I)
    \end{align*}
\end{lemma}


\begin{proof}
    Let $\vx_0\sim \mathcal{N}(\mu,\vect \Sigma)$. We have that 
    \begin{align*}
        \vx_t | \vx_0 &\sim \mathcal{N}(\sqrt{\alpha_t} \vx_0, (1-\alpha_t) \vect I)
        \\
        \Longrightarrow \vx_t &= \sqrt{\alpha_t}\vx_0 + \sqrt{1-\alpha_t} \vect z
    \end{align*}
    where $\vect z\sim \mathcal{N}(0,\vect I)$ is independent of $\vx_0$. Therefore,
    \begin{align*}
        \vx_t &\sim \mathcal{N}\fitparenth{\sqrt{\alpha_t} \vect \mu, \alpha_t \vect \Sigma +(1-\alpha_t)\vect I}
    \end{align*}
\end{proof}



As per \eqref{eq:ddpm-perturbation} $\alpha_t$ decreases from $1$ to $0$, the component means contract toward the origin and the covariances inflate toward $\mathbf{I}$. Substituting \eqref{eq:ddpm-perturbation} into \eqref{eq:responsibility-weights} reveals two limits of interest. At small $t$ ($\alpha_t \to 1$), the perturbed components remain well separated, and the responsibilities concentrate sharply on the closest mode, recovering the M-step regime of the EM algorithm. At large $t$ ($\alpha_t \to 0$), the perturbed components all approach $\mathcal{N}(\mathbf{0}, \mathbf{I})$ and the responsibilities flatten toward the prior weights $(1-\gamma, \gamma)$, becoming nearly independent of ${\vx}$. Sensitivity of the score to the mixture parameter $\gamma$ is therefore concentrated at the noise scales for which the responsibility weights are neither saturated nor uninformative --- a window whose location and width are determined by the geometry of the components.

%\subsubsection{Learning Diffusion score matching }
We now demonstrate an equivalence between the neural network loss \eqref{eqn: neural network loss} and the diffusion score matching loss \eqref{eq:diffScoreMatchingLoss}. 
\begin{theorem}
\label{lem:nn-dsm-ordering}
Let $\pdist{\gamma}$ be a data distribution on $\mathbb{R}^d$ with $\mathbb{E}_{\pdist{\gamma}}\|\vx_0\|^2 < \infty$, and let $\{\pdist[t]{\gamma}\}_{t\in[0,1]}$ be its marginals under the VP-SDE \eqref{eq:forward-sde}--\eqref{eq:forward-kernel}. Let $\vs_\btheta:\mathbb{R}^d \times [0,1] \to \mathbb{R}^d$ be a score model with $\mathbb{E}_{\pdist[t]{\gamma}}\|\vs_\btheta(\cdot,t)\|^2 < \infty$ for almost every $t \in [0,1]$. Then
\begin{equation}
   \lnn(\pdist{\btheta}, \pdist{\gamma}) \;=\; \ldsm(\pdist{\btheta}, \pdist{\gamma}) \;+\; C(\pdist{\gamma}),
   \label{eq:nn-dsm-equivalence}
\end{equation}
where the constant $C(\pdist{\gamma})$ does not depend on $\btheta$ and admits the closed form
\begin{align}   \label{eq:nn-dsm-constant}
   C(\pdist{\gamma}) &\;=\; \int_0^1 \lambda_t \cdot \frac{\alpha_t}{(1-\alpha_t)^2}\,\mathrm{MMSE}_t(\pdist{\gamma})\, dt,\\
   %\qquad
\mathrm{MMSE}_t(\pdist{\gamma}) 
&\;:=\; 
\mathbb{E}_{\substack{\vx_0 \,\sim\, \pdist{\gamma} \\ \vx_t \,\sim\, \pdist[t]{\gamma}(\cdot\mid\vx_0)}}\!\left[\,\big\|\vx_0 - \widehat{\vx}_0(\vx_t)\big\|^2\,\right],\\
\widehat{\vx}_0(\vx_t) &\;:=\; \mathbb{E}_{\vx_0' \,\sim\, \pdist{\gamma}(\cdot\mid\vx_t)}[\vx_0'].
%\mathrm{MMSE}_t(\pdist{\gamma}) := \mathbb{E}\big\|\vx_0 - \mathbb{E}[\vx_0 \mid \vx_t]\big\|^2.
\end{align}
Moreover, $C(\pdist{\gamma}) \ge 0$, with equality if and only if $\pdist{\gamma}$ is a Dirac mass. In particular,
\begin{equation}
   \lnn(\pdist{\btheta}, \pdist{\gamma}) \;\ge\; \ldsm(\pdist{\btheta}, \pdist{\gamma})
   \label{eq:nn-dsm-ordering}
\end{equation}
for every $\btheta \in \boldsymbol{\Theta}$.
\end{theorem}


\begin{proof}
Our arguments follow the reasoning in \cite{6795935} with an additional logic for guaranteeing the required relation between the neural network loss and the diffusion score matching loss.  
%[Proof of \cref{lem:nn-dsm-ordering}]
Throughout, all expectations are taken under $\pdist{\gamma}(\vx_0)$, $\pdist[t]{\gamma}(\vx_t\mid \vx_0)$ for the joint, and under $\pdist[t]{\gamma}(\vx_t)$ for the marginal. We suppress the dependence on $t$ in $\vs_\btheta(\vx_t, t)$ when convenient.
%\paragraph{Step 1: Pointwise (per-$t$) decomposition.} 
Fix $t \in [0,1]$. Expanding the squared norm in the inner expectation of $\ldsm$,
\begin{equation}
   \mathbb{E}_{\vx_t}\!\big\|\vs_\btheta(\vx_t,t) - \score[t]{\gamma}(\vx_t)\big\|^2
   = \mathbb{E}_{\vx_t}\|\vs_\btheta(\vx_t,t)\|^2 \,-\, 2\,b_t(\btheta) \,+\, \mathbb{E}_{\vx_t}\!\big\|\score[t]{\gamma}(\vx_t)\big\|^2,
   \label{eq:proof-esm-expand}
\end{equation}
with $b_t(\btheta) := \mathbb{E}_{\vx_t}\!\big\langle \vs_\btheta(\vx_t,t),\, \score[t]{\gamma}(\vx_t)\big\rangle$. Similarly,
\begin{align}\label{eq:proof-dsm-expand}
&\mathbb{E}_{\vx_0,\vx_t}\!\big\|\vs_\btheta(\vx_t,t) - \nabla_{\vx_t}\log \pdist[t]{\gamma}(\vx_t \mid \vx_0)\big\|^2\\\nonumber 
   &\qquad = \mathbb{E}_{\vx_0,\vx_t}\|\vs_\btheta(\vx_t, t)\|^2 \,-\, 2\,a_t(\btheta) \,+\, \mathbb{E}_{\vx_0,\vx_t}\!\big\|\nabla_{\vx_t}\log \pdist[t]{\gamma}(\vx_t\mid\vx_0)\big\|^2,
\end{align}
with $a_t(\btheta) := \mathbb{E}_{\vx_0,\vx_t}\!\big\langle \vs_\btheta(\vx_t,t),\, \nabla_{\vx_t}\log \pdist[t]{\gamma}(\vx_t\mid\vx_0)\big\rangle$. 
The first term in each expansion coincides because $\vs_\btheta(\vx_t,t)$ depends only on $\vx_t$. We now claim that for all $\theta$, $b_t(\btheta) = a_t(\btheta)$. Using $\score[t]{\gamma} = \nabla \pdist[t]{\gamma}/\pdist[t]{\gamma}$ and the marginalization $\pdist[t]{\gamma}(\vx_t) = \int \pdist[t]{\gamma}(\vx_t\mid\vx_0)\,\pdist{\gamma}(\vx_0)\, d\vx_0$,
\begin{align*}
   b_t(\btheta) 
   &= \int \pdist[t]{\gamma}(\vx_t)\,\big\langle \vs_\btheta(\vx_t,t),\, \score[t]{\gamma}(\vx_t)
   %\nabla_{\vx_t}\log \pdist[t]{\gamma}(\vx_t)
   \big\rangle\, d\vx_t \\
   &= \int \big\langle \vs_\btheta(\vx_t,t),\, \nabla_{\vx_t}\pdist[t]{\gamma}(\vx_t)\big\rangle\, d\vx_t \\
   &= \int \Big\langle \vs_\btheta(\vx_t,t),\, \nabla_{\vx_t}\!\!\int \pdist[t]{\gamma}(\vx_t\mid\vx_0)\,\pdist{\gamma}(\vx_0)\, d\vx_0\Big\rangle\, d\vx_t \\
   &= \iint \big\langle \vs_\btheta(\vx_t,t),\, \nabla_{\vx_t}\pdist[t]{\gamma}(\vx_t\mid\vx_0)\big\rangle\,\pdist{\gamma}(\vx_0)\, d\vx_0\, d\vx_t \\
   &= \iint \pdist[t]{\gamma}(\vx_t\mid\vx_0)\,\pdist{\gamma}(\vx_0)\,\big\langle \vs_\btheta(\vx_t,t),\, \nabla_{\vx_t}\log \pdist[t]{\gamma}(\vx_t\mid\vx_0)\big\rangle\, d\vx_0\, d\vx_t \\
   &= a_t(\btheta),
\end{align*}
where Leibniz's rule applies by smoothness and Gaussian decay of the forward kernel, and Fubini's theorem applies due to the regularity assumption on $\vs_{\btheta}$. Combining the above,
\begin{equation}
   \mathbb{E}_{\vx_0,\vx_t}\!\big\|\vs_\btheta(\vx_t, t) - \nabla_{\vx_t}\log \pdist[t]{\gamma}(\vx_t\mid\vx_0)\big\|^2
   - \mathbb{E}_{\vx_t}\!\big\|\vs_\btheta(\vx_t, t) - \score[t]{\gamma}(\vx_t) \big\|^2
   = \mathrm{gap}(t),
   \label{eq:proof-residual}
\end{equation}
where
\begin{equation}
   \mathrm{gap}(t) := \mathbb{E}_{\vx_0,\vx_t}\!\big\|\nabla_{\vx_t}\log \pdist[t]{\gamma}(\vx_t\mid\vx_0)\big\|^2 - \mathbb{E}_{\vx_t}\!\big\|\score[t]{\gamma}(\vx_t)\big\|^2,
   \label{eq:proof-Dt-def}
\end{equation}
which is independent of $\btheta$. The Gaussian forward kernel \eqref{eq:forward-kernel} gives
\begin{equation}
   \nabla_{\vx_t}\log \pdist[t]{\gamma}(\vx_t\mid\vx_0)
   = -\frac{\vx_t - \sqrt{\alpha_t}\vx_0}{1-\alpha_t}
   = -\frac{\vect{z}}{\sqrt{1-\alpha_t}},\quad \vect{z}\sim\mathcal{N}(\vect{0},\vect{I}),
   \label{eq:proof-cond-score}
\end{equation}
hence the first term of \eqref{eq:proof-Dt-def} equals
\begin{equation}
\mathbb{E}~\!\big\|\nabla_{\vx_t}\log \pdist[t]{\gamma}(\vx_t\mid\vx_0)\big\|^2 = \frac{d}{1-\alpha_t}.
   \label{eq:proof-cond-norm}
\end{equation}
For the marginal score, take expectations of \eqref{eq:proof-cond-score} conditional on $\vx_t$ we obtain Tweedie's formula:
\begin{equation}
   \score[t]{\gamma}(\vx_t)
   = \mathbb{E}~\!\big[\nabla_{\vx_t}\log \pdist[t]{\gamma}(\vx_t\mid\vx_0)\,\big|\,\vx_t\big]
   = \frac{\sqrt{\alpha_t}\,\mathbb{E}[\vx_0\mid\vx_t] - \vx_t}{1-\alpha_t}.
   \label{eq:proof-tweedie}
\end{equation}
Setting $\hat{\bm{\eta}}(\vx_t) := \sqrt{\alpha_t}\,\mathbb{E}[\vx_0\mid\vx_t] = \mathbb{E}[\sqrt{\alpha_t}\vx_0 \mid \vx_t]$, we see that, %. Then
\begin{equation}
   \mathbb{E}~\!\big\|\score[t]{\gamma}(\vx_t)\big\|^2 = \frac{\mathbb{E}\|\vx_t - \hat{\bm{\eta}}(\vx_t)\|^2}{(1-\alpha_t)^2}.
   \label{eq:proof-score-norm}
\end{equation}
Decomposing $\vx_t - \hat{\bm{\eta}} = (\vx_t - \sqrt{\alpha_t}\vx_0) + (\sqrt{\alpha_t}\vx_0 - \hat{\bm{\eta}}) = \sqrt{1-\alpha_t}\,\vect{z} + \sqrt{\alpha_t}(\vx_0 - \mathbb{E}[\vx_0\mid\vx_t])$ and expanding the expectation we get,
\begin{equation}
   \mathbb{E}\|\vx_t - \hat{\bm{\eta}}\|^2 = (1-\alpha_t)\,d + 2\sqrt{\alpha_t(1-\alpha_t)}\,\mathbb{E}~\!\big\langle\vect{z},\,\vx_0 - \mathbb{E}[\vx_0\mid\vx_t]\big\rangle + \alpha_t\,\mathrm{MMSE}_t(\pdist{\gamma}).
   \label{eq:proof-decomp}
\end{equation}
For the cross term, use $\vect{z} = (\vx_t - \sqrt{\alpha_t}\vx_0)/\sqrt{1-\alpha_t}$ and orthogonality of conditional expectation, $\mathbb{E}~\!\big[(\vx_0 - \mathbb{E}[\vx_0\mid\vx_t])\,\widehat{\vx}_0(\vx_t)\big] = 0$ for any $g$:
\begin{align*}
   \mathbb{E}~\!\big\langle\vect{z},\,\vx_0 - \mathbb{E}[\vx_0\mid\vx_t]\big\rangle
   &= \tfrac{1}{\sqrt{1-\alpha_t}}\,\mathbb{E}~\!\big\langle \vx_t - \sqrt{\alpha_t}\vx_0,\,\vx_0 - \mathbb{E}[\vx_0\mid\vx_t]\big\rangle \\
   &= \tfrac{1}{\sqrt{1-\alpha_t}}\,%\bi 
   \widehat{\vx}_0(0 - \sqrt{\alpha_t}\,\mathrm{MMSE}_t(\pdist{\gamma})\big) \\
   &= -\frac{\sqrt{\alpha_t}\,\mathrm{MMSE}_t(\pdist{\gamma})}{\sqrt{1-\alpha_t}}.
\end{align*}
Substituting into \eqref{eq:proof-decomp},
\begin{equation}
   \mathbb{E}\|\vx_t - \hat{\bm{\eta}}\|^2 = (1-\alpha_t)\,d - \alpha_t\,\mathrm{MMSE}_t(\pdist{\gamma}),
\end{equation}
and therefore via \eqref{eq:proof-score-norm},
\begin{equation}
   \mathbb{E}~\!\big\|\score[t]{\gamma}(\vx_t)\big\|^2 = \frac{d}{1-\alpha_t} - \frac{\alpha_t\,\mathrm{MMSE}_t(\pdist{\gamma})}{(1-\alpha_t)^2}.
   \label{eq:proof-score-final}
\end{equation}
Combining \eqref{eq:proof-cond-norm} and \eqref{eq:proof-score-final} in \eqref{eq:proof-Dt-def},
\begin{equation}
   \mathrm{gap}(t) = \frac{\alpha_t\,\mathrm{MMSE}_t(\pdist{\gamma})}{(1-\alpha_t)^2} \;\ge\; 0.
   \label{eq:proof-Dt-final}
\end{equation}
Hence, multiplying \eqref{eq:proof-residual} by $\lambda_t$ and integrating over $t\sim \mathrm{Unif}(0,1)$, the difference in the neural network loss and the diffusion score matching loss can be written as, 
\begin{equation}
   \lnn(\pdist{\btheta}, \pdist{\gamma}) - \ldsm(\pdist{\btheta}, \pdist{\gamma})
   = \int_0^1 \lambda_t\,\mathrm{gap}(t)\,dt
   = \int_0^1 \lambda_t\,\frac{\alpha_t\,\mathrm{MMSE}_t(\pdist{\gamma})}{(1-\alpha_t)^2}\,dt
   =: C(\pdist{\gamma}),
\end{equation}
proving \eqref{eq:nn-dsm-equivalence}--\eqref{eq:nn-dsm-constant}.
For the equality case, note that $C(\pdist{\gamma}) = 0$ iff $\mathrm{MMSE}_t(\pdist{\gamma}) = 0$ for almost every $t \in (0,1)$. The latter requires $\vx_0 = \mathbb{E}[\vx_0\mid\vx_t]$ almost surely, i.e., $\vx_0$ is $\sigma(\vx_t)$-measurable. Since $\vx_t = \sqrt{\alpha_t}\vx_0 + \sqrt{1-\alpha_t}\vect{z}$ with $\vect{z}$ independent of $\vx_0$ and $1-\alpha_t > 0$, this forces $\vx_0$ to be deterministic.
\end{proof}




\subsection{Proof of  \Cref{thm: GMM Sensitivity Emerges}}\label{Proof of Theorem 1}



\begin{remark}
    In this theorem, we use $\lambda_t=1$ rather than $\lambda_t= 1-\alpha_t$. 
    The value of $t$ where $\lsm(\pdist[t]{\gamma},\pdist[t]{\gamma^*})$ is large correspond to when $\alpha_t$ is close to $0$ for when the modes are reasonably well-separated, so that $1-\alpha_t$ is close to $1$. The same analysis can be performed in the presence of any weighting function $\lambda_t$ if one is willing to accept a more complicated expression for the constant $C$; we elaborate on this in \Cref{prop: Lower bound integral}. For all of our numerical results, we use the convention that $\lambda_t=1-\alpha_t$ to show that our results still hold for the weighting factor used in practice.
\end{remark}


Throughout the proof, let $\separation:=\norm{\mu_1-\mu_0}$. From \Cref{lemma:mixture-closure} and \Cref{lemma:mix_over_time}, we can generate samples from $\pdist[t]{\gamma}$ by sampling $c\sim \mathrm{Ber}(\gamma)$ and sampling $\vect y_t\sim \pdist[t]{c}$. Furthermore, the next lemma will show that we can generate these samples using the Probability Flow ODE (\citep{song2021scorebasedgenerativemodelingstochastic}) instead of the stochastic forward process.

First, note that 
\begin{align}\label{eq: mode decomposition}
    \ldsm(\pdist{\hat\gamma},\pdist{\gamma^*})&= \expect[t,\vect x_t\sim \pdist[t]{\gamma^*}]{\norm{\score[t]{\gamma^*}(\vect x_t)-\score[t]{\hat\gamma}}(\vect x_t)}
    \\
    &= (1-\gamma^*)\expect[t,\vect x_t\sim \pdist[t]{\gamma^*}|c=0]{\norm{\score[t]{\gamma^*}(\vect x_t)-\score[t]{\hat\gamma}}(\vect x_t)}
    \\
    &\quad\quad + \gamma^*\expect[t,\vect x_t\sim \pdist[t]{\gamma^*}|c=1]{\norm{\score[t]{\gamma^*}(\vect x_t)-\score[t]{\hat\gamma}}(\vect x_t)}
    \\
    &= (1-\gamma^*)\expect[t,\vect x_t\sim \pdist[t]{0}]{\norm{\score[t]{\gamma^*}(\vect x_t)-\score[t]{\hat\gamma}}(\vect x_t)}
    \\
    &\quad\quad + \gamma^*\expect[t,\vect x_t\sim \pdist[t]{1}]{\norm{\score[t]{\gamma^*}(\vect x_t)-\score[t]{\hat\gamma}}(\vect x_t)}
\end{align}

In particular, we can lower-bound $\ldsm$ by independently lower-bounding each of these expectations.

\begin{lemma}
    Consider any distribution $\pi$ and let $\pi_t$ be the noisy marginal at time $t$ of the forward process. Then, if $\vect y_t$ satisfies the Probability Flow ODE
    \begin{equation}\label{eqn: Probability Flow ODE}
        \mathrm{d}\vect y_t = -\frac{1}{2}\beta_t\fitparenth{\vect y_t + \nabla \log p_{\vect y_t}(\vect y_t)}\mathrm{d}t
    \end{equation}
    and $\vect y_0\sim \pi$, then $\vect y_t\sim \pi_t$.
\end{lemma}

\begin{proof}
    Let $\vx_t$ satisfy the Variance Preserving SDE and $\vect y_t$ be as in the lemma statement. Then, we have the Fokker-Planck equations
    \begin{align*}
        \frac{\partial p_{\vx_t} }{\partial t}&= -\nabla \cdot \fitparenth{-\frac{1}{2}\beta_t \vx p_{\vx_t} } + \frac{1}{2}\beta_t\Delta p_{\vx_t} 
        \\
        &=\frac{1}{2}\beta_t  \nabla \cdot \fitparenth{\vx p_{\vx_t} } + \frac{1}{2}\beta_t\Delta p_{\vx_t} 
        \\
        \\
        \frac{\partial p_{\vect y_t} }{\partial t}
        &= \frac{1}{2}\beta_t \nabla \cdot \fitparenth{\vect y p_{\vect y_t} }  +\frac{1}{2}\beta_t \nabla \cdot \fitparenth{p_{\vect y_t}\nabla \log p_{\vect y_t} } 
        \\
        &= \frac{1}{2}\beta_t \nabla \cdot \fitparenth{\vect y p_{\vect y_t} }  +\frac{1}{2}\beta_t \Delta p_{\vect y_t}
    \end{align*}
\end{proof}

The Probability Flow ODE (\eqref{eqn: Probability Flow ODE}) substantially simplifies the analysis because of the following fact:
\begin{lemma}\label{lemma: simple paths}
    For $\vect y_0$, and $\vect y_t$ following the Probability Flow ODE associated with $c$ and starting from $\vect y_0\sim \pdist{c}$, we have that 
    \[\vect y_t - \sqrt{\alpha_t}\mu_c = \vect y_0 - \mu_c\]
\end{lemma}

\begin{proof}
    Note that $\frac{\mathrm d}{\mathrm d t}\sqrt{\alpha_t}=\frac{1}{2}\beta_t \sqrt{\alpha_t}$. Therefore,
    \begin{align*}
        \frac{\mathrm d}{\mathrm dt}(\vect y_t -\sqrt{\alpha_t}\mu_c) &=-\frac{1}{2}\beta_t \fitparenth{\vect y_t + \nabla \log \pdist[t]{c}(\vect y_t)} - \frac{1}{2}\beta_t \sqrt{\alpha_t} \mu_c
        \\
        &=-\frac{1}{2}\beta_t \fitparenth{\vect y_t + \sqrt{\alpha_t} \mu_c - \vect y_t - \sqrt{\alpha_t} \mu_c}
        \\
        &=0
    \end{align*}
\end{proof}

We now focus on the case where $c=0$ (the case of $c=1$ will follow via symmetry). Let $\vect z_t = \vect y_t|c=0$. Our first goal is to lower bound $\norm{\score[t]{\gamma^*}(\vect z_t) - \score[t]{\hat \gamma}(\vect z_t)}^2$ pointwise in $t$ and $\vect z_t$. We first show an analytic expression for this score difference, then identify a range of $t$'s along each path where this expression is sufficiently large. Finally, by integrating over initial conditions $\vect z_0$, we get a lower bound for $\expect[t,\vect z_t\sim \pdist[t]{0}]{\norm{\score[t]{\gamma^*}(\vect z_t)-\score[t]{\hat\gamma}(\vect z_t)}^2}$.
\begin{proposition}\label{prop:exact norm expr}
    Let $\hat \gamma, \gamma^*\in [0,1]$. Then, 
    \[\norm{\score[t]{\hat \gamma}(\vect z_t) - \score[t]{\gamma^{*}}(\vect z_t)}^2 = ({\hat\gamma}-\gamma^*)^2  \frac{\norm{\mu_0-\mu_1}^2 \alpha_te^{2f_0(\vect z_0, t)}}{(1-{\hat\gamma} + {\hat\gamma}e^{f_0(\vect z_0, t)})^2 (1-\gamma^* + \gamma^* e^{f_0(\vect z_0, t)})^2}\]
    where 
    \[f_0(\vect z_0, t):=-\frac{1}{2}\norm{\mu_1-\mu_0}^2\alpha_t + \inprod{\vect z_0 - \mu_0, \mu_1 - \mu_0}\sqrt{\alpha_t}\]
\end{proposition}

\begin{proof}
    We \[\pdist[t]{0}(\vx)=\frac{1}{Z_t}\exp\fitparenth{-\frac{1}{2}\norm{\vx-\sqrt{\alpha_t}\mu_0}^2}\]
     \[\pdist[t]{1}(\vx)=\frac{1}{Z_t}\exp\fitparenth{-\frac{1}{2}\norm{\vx-\sqrt{\alpha_t}\mu_1}^2}\]

    Expanding the log likelihoods, we have
    \begin{align*}
        \log \frac{p_1(\vect z_t)}{p_0(\vect z_t)}&= -\frac{1}{2} \norm{\mu_0-\mu_1}^2 \alpha_t + \inprod{\vect z_0 - \mu_0, \mu_1 - \mu_0}\sqrt{\alpha_t}
    \end{align*}

    Let $f(\vect z_0,t) :=\log \frac{p_1(\vect z_t)}{p_0(\vect z_t)}$. Then,
    \begin{align*}
        \score[t]{\gamma}(\vect z_t)&= \frac{(1-\gamma)\;\pdist[t]{0}(\vect z_t)\score[t]{0}(\vect z_t) + \gamma\;\pdist[t]{1}(\vect z_t)\score[t]{1}(\vect z_t)}{\pdist[t]{\gamma}(\vect z_t)}
        \\
        &=\score[t]{0}(\vect z_t)  +  \frac{  \gamma\;\pdist[t]{1}(\vect z_t)\fitparenth{\score[t]{1}(\vect z_t)-\score[t]{0}(\vect z_t)}}{\pdist[t]{\gamma}(\vect z_t)}
        \\
        &=\score[t]{0}(\vect z_t)  +  \frac{  \gamma\;\pdist[t]{1}(\vect z_t)}{\pdist[t]{\gamma}(\vect z_t)}\sqrt{\alpha_t}\fitparenth{ \mu_1- \mu_0}
        \\
        &=\score[t]{0}(\vect z_t)  +  \frac{  \gamma\;e^{f(\vect z_0,t)}}{1-\gamma+\gamma \; e^{f(\vect z_0,t)}}\sqrt{\alpha_t}\fitparenth{ \mu_1- \mu_0}
    \end{align*}
    where we used $\nabla \log \pdist[t]{c}(\vx) = \sqrt{\alpha_t}\mu_c-\vx$. Therefore,
    \begin{align*}
        &\score[t]{\hat \gamma}(\vect z_t) - \score[t]{\gamma^{*}}(\vect z_t)= \frac{  ({\hat\gamma}-\gamma^*)\;e^{f(\vect z_0,t)}}{\fitparenth{1-{\hat\gamma}+{\hat\gamma} \; e^{f(\vect z_0,t)}}\fitparenth{1-\gamma^*+\gamma^* \; e^{f(\vect z_0,t)}}}\sqrt{\alpha_t}\fitparenth{ \mu_1- \mu_0}
    \end{align*}

    This gives
    \[\norm{\score[t]{\hat \gamma}(\vect z_t) -\score[t]{\gamma^{*}}(\vect z_t)}^2=({\hat\gamma}-\gamma^*)^2\frac{ \norm{ \mu_1- \mu_0}^2\alpha_t\;e^{2f(\vect z_0,t)}}{\fitparenth{1-{\hat\gamma}+{\hat\gamma} \; e^{f(\vect z_0,t)}}^2\fitparenth{1-\gamma^*+\gamma^* \; e^{f(\vect z_0,t)}}^2} \]
\end{proof}

\begin{lemma}
    For all $x\in \RR$ and $\gamma\in [0,1]$, we have 
    \[\frac{e^x}{(1-\gamma + \gamma e^x)^2}\geq \frac{1}{4\gamma (1-\gamma)}e^{-\frac{1}{4}\fitparenth{x-\log \frac{1-\gamma}{\gamma}}^2}\]
    with equality only when $x=\log \frac{1-\gamma}{\gamma}$. In particular, 
    \begin{align*}
        &\frac{e^{2f_0(\vect z_0, t)}}{\fitparenth{1-{\hat\gamma} + {\hat\gamma} e^{f_0(\vect z_0, t)}}^2 \fitparenth{1-\gamma^* + \gamma^* e^{f_0(\vect z_0, t)}}^2}\geq \frac{e^{-\frac{1}{4}\fitparenth{f_0(\vect z_0,t)-\log \frac{1-{\hat\gamma}}{{\hat\gamma}}}^2-\frac{1}{4}\fitparenth{f_0(\vect z_0,t)-\log \frac{1-\gamma^*}{\gamma^*}}^2}}{16 {\hat\gamma} (1-{\hat\gamma}) \gamma^* (1-\gamma^*)}
    \end{align*}
\end{lemma}

\begin{proof}
    Let $g(x)= \frac{e^x}{(1-\gamma + \gamma e^x)^2}$ and $h(x)=\frac{1}{4\gamma (1-\gamma)}e^{-\frac{1}{4}\fitparenth{x-\log \frac{1-\gamma}{\gamma}}^2}$. Then, 
    \begin{align*}
        \frac{\mathrm d}{\mathrm d x}\log g(x) &=\frac{1-\gamma -\gamma e^x}{1-\gamma + \gamma e^x}
    \end{align*}

    For all $\gamma\in [0,1]$ $1-\gamma + \gamma e^x>0$. Therefore, $x^*:=\log \frac{1-\gamma}{\gamma}$ is the only critical point of $g(x)$, and $g(x^*)=\frac{1}{4\gamma\fitparenth{1-\gamma}}$. Similarly $h(x^*)=\frac{1}{4\gamma\fitparenth{1-\gamma}}$ and $\frac{\mathrm d}{\mathrm d x}\log h(x^*)=0$. We have
    \begin{align*}
        \frac{\mathrm{d}}{\mathrm{d}x}[\log h(x)-\log g(x)]
        &=-\frac{1}{2}(x-x^*) - \frac{1-e^{x-x^*}}{1+e^{x-x^*}}
        \\
        \frac{\mathrm{d}^2}{\mathrm{d}x^2}[\log h(x)-\log g(x)]
        &= -\frac{(1-e^{x-x^*})^2}{2(1+e^{x-x^*})^2}
    \end{align*}
    
    So, $\frac{\mathrm{d}^2}{\mathrm{d}x^2}[\log h(x)-\log g(x)]\leq 0$ with equality only when $x=x^*$. Therefore, $\log h-\log g$ is strictly concave away from $x^*$ and maximized at $x^*$, where $h(x^*)=g(x^*)$. Therefore, $h< g$.
\end{proof}

\begin{lemma}\label{lemma:bounds}
    Let $\vect z_t$ and $\vect z_0$ be as above. Let $M(\vect z_0):=\inprod{\vect z_0 - \mu_0, \mu_1 - \mu_0}$, and define the log-odds ratios 
    \begin{align*}
        \hat r &:= \log \frac{1-{\hat\gamma}}{{\hat\gamma}}
        \\
        r^* &:= \log \frac{1-\gamma^*}{\gamma^*}
    \end{align*}
    
    Let $\epsilon>0$ be such that 
    \begin{align*}
    \epsilon &\leq \exp\fitparenth{-\frac{1}{8}\fitbracket{(\hat r-r^*)^2 + \fitparenth{\hat r+r^*-\frac{M(\vect z_0)^2}{\separation^2}}_+^2}}
    \end{align*}
    where $(\cdot)_+:= \max\set{\cdot, 0}$.
    
    Let
    \begin{align*}
        a&=\max\left\{\frac{M(\vect z_0)}{\separation^2} + \frac{1}{\separation}\sqrt{\fitparenth{\frac{M(\vect z_0)^2}{\separation^2} - \fitparenth{ \hat r + r^*}-\sqrt{8\log \frac{1}{\epsilon}-\fitparenth{\hat r-r^*}^2}}_+},\sqrt{\alpha_1}\right\}
        \\
        b&= \min\left\{\frac{M(\vect z_0)}{\separation^2} + \frac{1}{\separation}\sqrt{\frac{M(\vect z_0)^2}{\separation^2} - \fitparenth{ \hat r + r^*}+\sqrt{8\log \frac{1}{\epsilon}-\fitparenth{\hat r-r^*}^2}},1\right\}
    \end{align*}

    Then for all $t$ such that $\sqrt{\alpha_t}\in [a,b]$, 
    \[e^{-\frac{1}{4}\fitparenth{f_0(\vect z_0,t)-\hat r }^s2-\frac{1}{4}\fitparenth{f_0(\vect z_0,t)-r^*}^2} \geq \epsilon\]
\end{lemma}

\begin{proof}
    First, the inequality is equivalent to 
    \begin{align*}
       \fitparenth{f_0(\vect z_0,t)-r^*}^2+\fitparenth{f_0(z_0,t)-r^*}^2\leq 4\log \frac{1}{\epsilon}
    \end{align*}

    
    
    
    Plugging in the definition of $f$, the inequality holds whenever 
    \[\fitbracket{\frac{\separation^2}{2} \fitparenth{\sqrt{\alpha_t}-\frac{M(\vect z_0)}{\separation^2}}^2 - \frac{M(\vect z_0)^2}{2\separation^2}+ \frac{\hat r+r^*}{2}}^2 \leq 2\log \frac{1}{\epsilon}-\frac{(\hat r-r^*)^2}{4} \]
    
    We wish to identify the values of $\sqrt{\alpha_t}$ which make this inequality hold. This function is continuous, so its sign can only change at the roots of the corresponding equality constraint:
    \[\fitbracket{\frac{\separation^2}{2} \fitparenth{\sqrt{\alpha_t}-\frac{M(\vect z_0)}{\separation^2}}^2 - \frac{M(\vect z_0)^2}{2\separation^2}+ \frac{\hat r+r^*}{2}}^2 = 2\log \frac{1}{\epsilon}-\frac{(\hat r -r^*)^2}{4}\]

    Solving this quartic gives the roots
    \[\sqrt{\alpha_t}=\frac{M(\vect z_0)}{\separation^2} \pm\sqrt{\frac{M(\vect z_0)^2}{\separation^4}-\frac{\hat r+r^*}{\separation^2}\pm \frac{1}{\separation^2}\sqrt{8\log \frac{1}{\epsilon}-(\hat r-r^*)^2}}\]
    
    The condition on $\epsilon$ is such that the square roots are all well-defined. Let $b'$ be the largest of these roots (both $\pm$'s are set to $+$). As $\sqrt{\alpha_t}\rightarrow \infty$, we would have $f_0(\vect z_0,t)\rightarrow -\infty$ (because $-\separation^2/2$ is strictly negative), and hence 
    \[e^{-\frac{1}{4}\fitparenth{f_0(\vect z_0,t)-\hat r }^2-\frac{1}{4}\fitparenth{f_0(\vect z_0,t)-r^*}^2} \rightarrow 0\]
    
    Therefore, for $\sqrt{\alpha_t}>b'$ the inequality does not hold. Note that the four roots above are distinct (with two being complex for $\epsilon$ sufficiently small), meaning that each root has multiplicity 1, and hence the polynomial changes signs at $b'$. Therefore, for $\sqrt{\alpha_t}$ between $b'$ and the next largest real root, the inequality holds. The next largest real root is
    \[a'=\begin{cases}
        \frac{M(\vect z_0)}{\separation^2} +\sqrt{\frac{M(\vect z_0)^2}{\separation^4}-\frac{\hat r+r^*}{\separation^2}-\frac{1}{\separation^2}\sqrt{8\log \frac{1}{\epsilon}-(\hat r-r^*)^2}} & \text{if}\quad \frac{M(\vect z_0)^2}{\separation^2}-(\hat r+r^*)-\sqrt{8\log \frac{1}{\epsilon}-(\hat r-r^*)^2} \geq0
        \\
        \frac{M(\vect z_0)}{\separation^2} -\sqrt{\frac{M(\vect z_0)^2}{\separation^4}-\frac{\hat r+r^*}{\separation^2}+\frac{1}{\separation^2}\sqrt{8\log \frac{1}{\epsilon}-(\hat r-r^*)^2}} & \text{otherwise}
    \end{cases}\]
    
    In the second case we have $a'\leq \frac{M(\vect z_0)}{\separation^2}$, so
    \[a'\leq \frac{M(\vect z_0)}{\separation^2} +\sqrt{\fitparenth{\frac{M(\vect z_0)^2}{\separation^4}-\frac{\hat r+r^*}{\separation^2}-\frac{1}{\separation^2}\sqrt{8\log \frac{1}{\epsilon}-(\hat r-r^*)^2}}_+}=:a'' \]
    
    Therefore, the inequality holds on $[a',b']$ and $[a'',b']\subseteq [a',b']$, so the inequality holds on $[a'',b']$. Note that for $t\in [0,1]$, $\alpha_t\in [\alpha_1, 1]$ where $\alpha_1 >0$ is small, so we let $b:= \min\set{b', 1}$ and $a:=\max\set{a,\sqrt{\alpha_1}}$.
\end{proof}


\begin{proposition}\label{prop: lower bound over t}
    For all $\epsilon$ satisfying the condition in  \Cref{lemma:bounds}, for all $t$ s.t. $\sqrt{\alpha_t}\in[a,b]$ (with $a$ and $b$ as defined in the previous lemma) we have
    \begin{align*}
        \norm{\score[t]{\hat \gamma}(\vect z_t) - \score[t]{\gamma^{*}}(\vect z_t)}^2\geq\frac{(\hat \gamma-\gamma^*)^2\norm{ \mu_1- \mu_0}^2 \alpha_t}{16 \hat \gamma (1-\hat \gamma) \gamma^* (1-\gamma^*)}\epsilon
    \end{align*}
    
\end{proposition}

\begin{proof}
    Combining the above two lemmas,
    \begin{align*}
        \norm{\score[t]{\hat \gamma}(\vect z_t) - \score[t]{\gamma^{*}}(\vect z_t)}^2 &\geq\frac{(\hat \gamma-\gamma^*)^2\norm{ \mu_1- \mu_0}^2 \alpha_t}{16 \hat \gamma (1-\hat \gamma) \gamma^* (1-\gamma^*)}e^{-\frac{1}{4}\fitparenth{f_0(\vect z_0,t)-\hat r}^2-\frac{1}{4}\fitparenth{f_0(\vect z_0,t)-r^*}^2}
        \\
        &\geq\frac{(\hat \gamma-\gamma^*)^2\norm{ \mu_1- \mu_0}^2 \alpha_t}{16 \hat \gamma (1-\hat \gamma) \gamma^* (1-\gamma^*)}\epsilon
    \end{align*}
\end{proof}

We now turn our attention to lower bounding 
\[\int_0^1\norm{\score[t]{\hat \gamma}(\vect z_t) - \score[t]{\gamma^{*}}(\vect z_t)}^2\mathrm{d}t\]
along a given path starting from $\vect {z _ 0} $. From the lemmas above, we have
\[\int_0^1  \norm{\nabla \log \pdist[t]{\hat \gamma}(\vect z_t) - \nabla \log \pdist[t]{\gamma^{*}}(\vect z_t)}^2\mathrm{d}t > \frac{({\hat \gamma}-\gamma^*)^2\norm{ \mu_1- \mu_0}^2}{16{\hat \gamma}(1-{\hat \gamma}) \gamma^* (1-\gamma^*)}\epsilon\int_{t_{b^2}}^{t_{a^2}}\alpha_t\mathrm{d}t\] 
where $t_{a^2}$ and $t_{b^2}$ are such that $\alpha_{t_{a^2}}={a^2}$ and $\alpha_{t_{b^2}}={b^2}$. Note $\alpha_t$ is a decreasing function of $t$, hence why the bounds of the integral are reversed.  If we let $\alpha^{-1}(\tau)$ be the inverse function of $\alpha_t$, such that $\alpha_{\alpha^{-1}(\tau)} = \tau$, then
\[\int_{t_{b^2}}^{t_{a^2}}\alpha_t\mathrm{d}t=2\int_{a}^b \frac{\tau}{\beta_{\alpha^{-1}(\tau^2)}}\mathrm{d}\tau\]

In particular, if we consider the linear noise schedule 
\[\beta_t:=\beta_1 t\]
then
\[\int_{t_{b^2}}^{t_{a^2}}\alpha_t\mathrm{d}t=\frac{1}{\sqrt{\beta_1}}\int_{a}^b \frac{\tau}{\sqrt{-\ln \tau}}\mathrm{d}\tau\]

\begin{proposition}\label{prop: Lower bound integral}
    Let $a$, $b$, $\separation$, $z_0$, $\hat r$, $r^*$, and $M$, be as above, and assume $\epsilon$ satisfies the bound from \Cref{lemma:bounds}. Define 
    \[R_0:=\sqrt{8\ln \frac{1}{\epsilon}-(\hat r-r^*)^2}-(\hat r+r^*)\]
    
    Then, when $\frac{M(\vect z_0)}{\separation^2}\geq \sqrt{\alpha_1}$ and
    \[\epsilon \geq \exp\fitparenth{-\frac{1}{8}\fitbracket{(\separation^2 - 2M(\vect z_0)+(\hat r+r^*))^2 + (\hat r-r^*)^2}}\]
    then
    \begin{align*}
        \int_{a}^b \frac{\tau}{\sqrt{-\ln \tau}}\mathrm{d}\tau&\geq \frac{\frac{M(\vect z_0)}{\separation}\sqrt{\frac{M(\vect z_0)^2}{\separation^2} + R_0}+\frac{M(\vect z_0)^2}{2\separation^2} + \frac{1}{2}R_0}{\separation^2\sqrt{\ln(\separation)- \ln\fitparenth{\frac{M(\vect z_0)}{\separation}+ \frac{1}{2}\sqrt{\frac{M(\vect z_0)^2}{\separation^2} + R_0}}}}
    \end{align*}
\end{proposition}

\begin{remark}
    This proposition requires $\lambda_t=1$, as the analysis is far more challenging when the integrand $g$ is nonconvex. If one wishes to analyze alternative weighting factors or, indeed, an alternative noise schedule, all the analysis up to this point applies. The difference is that one now needs to lower bound the more challenging expression
    $\int_{t_{b^2}}^{t_{a^2}}\alpha_t\lambda_t\mathrm{d}t$. For the linear noise schedule and $\lambda_t=1-\alpha_t$, this take the form $\frac{1}{\sqrt{\beta_1}}\int_{a}^b \frac{\tau(1-\tau^2)}{\sqrt{-\ln \tau}}\mathrm{d}\tau$.
\end{remark}

\begin{proof}
    Let \[g(t):= \frac{\tau}{\sqrt{-\ln \tau}}\]

    Note that 
    \begin{align*}
        g''(t)&= \frac{1}{2t(-\ln t)^{3/2}} + \frac{3}{4t(-\ln t)^{5/2}} >0
    \end{align*}

    So, $g(t)$ is convex, meaning 
    \[\int_a^b g(\tau)\mathrm{d}\tau \geq (b-a)g\fitparenth{\frac{a+b}{2}}\]

    The condition that $\frac{M(\vect z_0)}{\separation^2}\geq \sqrt{\alpha_1}$ guarantees $a\geq\sqrt{\alpha_1}$. Similarly, the condition on $\epsilon$ guarantees $b\leq 1$, so
    \begin{align*}
        b-a&=\frac{1}{\separation}\sqrt{\frac{M(\vect z_0)^2}{\separation^2} - \fitparenth{ \hat r + r^*}+\sqrt{8\log \frac{1}{\epsilon}-\fitparenth{\hat r-r^*}^2}}
        \\
        &= \frac{1}{\separation}\sqrt{\frac{M(\vect z_0)^2}{\separation^2} + R_0}
    \end{align*}

    We have 
    \begin{align*}
        \frac{a+b}{2}&= \frac{M(\vect z_0)}{\separation^2}+\frac{1}{2\separation}\sqrt{\frac{M(\vect z_0)^2}{\separation^2} + R_0}
    \end{align*}
    and so
    \begin{align*}
        \Longrightarrow g\fitparenth{\frac{a+b}{2}}&= \frac{\frac{M(\vect z_0)}{\separation}+\frac{1}{2}\sqrt{\frac{M(\vect z_0)^2}{\separation^2} + R_0}}{\separation\sqrt{\ln(\separation)- \ln\fitparenth{\frac{M(\vect z_0)}{\separation}+ \frac{1}{2}\sqrt{\frac{M(\vect z_0)^2}{\separation^2} + R_0}}}}
    \end{align*}

    Combining these, 
    \begin{align*}
        \int_a^b g(\tau)\mathrm{d}\tau &\geq \frac{\frac{M(\vect z_0)}{\separation}\sqrt{\frac{M(\vect z_0)^2}{\separation^2} + R_0}+\frac{M(\vect z_0)^2}{2\separation^2} + \frac{1}{2}R_0}{\separation^2\sqrt{\ln(\separation)- \ln\fitparenth{\frac{M(\vect z_0)}{\separation}+ \frac{1}{2}\sqrt{\frac{M(\vect z_0)^2}{\separation^2} + R_0}}}}
    \end{align*}
\end{proof}



Combining this lower bound with what we showed before, we have 
\begin{align*}
    &\int_0^1\norm{\nabla \log \pdist[t]{\hat \gamma}(\vect z_t) - \nabla \log \pdist[t]{\gamma^{*}}(\vect z_t)}^2 \mathrm{d}t
    \\
    &\quad\quad\geq\frac{(\hat \gamma-\gamma^*)^2 \epsilon}{16\sqrt{\beta_1}\;\hat  \gamma (1-\hat \gamma) \gamma^* (1-\gamma^*)} \frac{\frac{M(\vect z_0)}{\separation}\sqrt{\frac{M(\vect z_0)^2}{\separation^2} + R_0}+\frac{M(\vect z_0)^2}{2\separation^2} + \frac{1}{2}R_0}{\sqrt{\ln(\separation)- \ln\fitparenth{\frac{M(\vect z_0)}{\separation}+ \frac{1}{2}\sqrt{\frac{M(\vect z_0)^2}{\separation^2} + R_0}}}}
\end{align*}

Recall that $\vect z_0\sim \mathcal{N}(\mu_0, \vect I)$, so 
\begin{align*}
    M(\vect z_0) &= \inprod{\vect z_0 -\mu_0, \mu_1-\mu_0} \sim \mathcal{N}(0, \norm{\mu_1-\mu_0}^2)
    \\
    \Longrightarrow \frac{M(\vect z_0)}{\separation}&\sim \mathcal{N}(0,1)
\end{align*}

\begin{proposition}\label{prop:c0_lowerbound}
    Let $\Phi$ be the CDF of $\mathcal{N}(0,1)$. Let $\tau\geq \separation \sqrt{\alpha_1}$, and let $\epsilon>0$ satisfy $ \epsilon \in (\epsilon_{\mathrm{min}}^0, \epsilon_{\mathrm{max}}^0)$ with 
    \begin{align*}
    \epsilon_{\mathrm{max}} ^0& =\exp\fitparenth{-\frac{1}{8}\fitbracket{(\hat r-r^*)^2 + \fitparenth{\hat r+r^*-\tau^2}_+^2}}
    \\
    \epsilon_{\mathrm{min}}^0 &= \exp\fitparenth{-\frac{1}{8}\fitbracket{(\separation(\separation - 2\tau)+(\hat r+r^*))^2 + (\hat r-r^*)^2}}
    \end{align*}
    
    Then, with probability $1-\Phi(\tau)$ over samples of $\vect z_0$, 
    \begin{align*}
        &\int_0^1\norm{\nabla \log \pdist[t]{\hat \gamma}(\vect z_t) - \nabla \log \pdist[t]{\gamma^{*}}(\vect z_t)}^2 \mathrm{d}t
        \\
        &\quad\quad\geq\frac{(\hat \gamma-\gamma^*)^2 \epsilon}{16\sqrt{\beta_1}\; \hat \gamma (1-\hat \gamma) \gamma^* (1-\gamma^*)} 
        \frac{\tau\sqrt{\tau^2 + R_0}+\frac{1}{2}\tau^2 + \frac{1}{2}R_0}{\sqrt{\ln(\separation)- \ln\fitparenth{\tau+ \frac{1}{2}\sqrt{\tau^2 + R_0}}}}
    \end{align*}
\end{proposition}

\begin{proof}
If we let $K:= \frac{M(\vect z_0)}{\separation}$ our bound becomes 
\begin{align*}
    &\int_0^1\norm{\nabla \log \pdist[t]{\hat \gamma}(\vect z_t) - \nabla \log \pdist[t]{\gamma^{*}}(\vect z_t)}^2 \mathrm{d}t
    \\
    &\quad\quad\geq\frac{(\hat\gamma-\gamma^*)^2 \epsilon}{16 \hat\gamma (1-\hat\gamma) \gamma^* (1-\gamma^*)} 
    \frac{K\sqrt{K^2 + R_0}+\frac{1}{2}K^2 + \frac{1}{2}R_0}{\sqrt{\ln(\separation)- \ln\fitparenth{K+ \frac{1}{2}\sqrt{K^2 + R_0}}}}
\end{align*}

Assume that $K\geq \tau$, which occurs with probability $1-\Phi(\tau)$.
Our bound is an increasing function in $K$ and is thus minimized at $K=\tau$. Similarly, our upper bound from before on $\epsilon$ is a decreasing function of $K$ and our lower bound from before is an increasing function of $\epsilon$, so our bounds on $\epsilon$ are tightest when we plug in $\tau$, giving $\epsilon_{\mathrm{max}}^0$ and $\epsilon_{\mathrm{min}}^0$ from the statement of the proposition.
\end{proof}

We now state the corresponding proposition for the case where $c=1$. Let $\vect w_t := \vect y_t |c=1$.
\begin{proposition}\label{prop:c1_lowerbound}
    Let $\tau\geq \separation \sqrt{\alpha_1}$, and let $\epsilon>0$ satisfy $ \epsilon \in (\epsilon_{\mathrm{min}}^1, \epsilon_{\mathrm{max}}^1)$ with 
    \begin{align*}
    \epsilon_{\mathrm{max}} ^1& =\exp\fitparenth{-\frac{1}{8}\fitbracket{(\hat r-r^*)^2 + \fitparenth{-(\hat r+r^*)-\tau^2}_+^2}}
    \\
    \epsilon_{\mathrm{min}}^1 &= \exp\fitparenth{-\frac{1}{8}\fitbracket{(\separation(\separation - 2\tau)+(\hat r+r^*))^2 - (\hat r-r^*)^2}}
    \end{align*}
    
    Then, with probability $1-\Phi(\tau)$ over samples of $\vect w_0$, 
    \begin{align*}
        &\int_0^1\norm{\nabla \log \pdist[t]{\hat \gamma}(\vect w_t) - \nabla \log \pdist[t]{\gamma^{*}}(\vect w_t)}^2 \mathrm{d}t
        \\
        &\quad\quad\geq\frac{(\hat \gamma-\gamma^*)^2 \epsilon}{16\sqrt{\beta_1}\; \hat \gamma (1-\hat \gamma) \gamma^* (1-\gamma^*)} 
        \frac{\tau\sqrt{\tau^2 + R_1}+\frac{1}{2}\tau^2 + \frac{1}{2}R_1}{\sqrt{\ln(\separation)- \ln\fitparenth{\tau+ \frac{1}{2}\sqrt{\tau^2 + R_1}}}}
    \end{align*}
    where
    \[R_1:=\sqrt{8\ln \frac{1}{\epsilon}-(\hat r-r^*)^2}+(\hat r+r^*)\]
\end{proposition}

\begin{proof}
    This follows from repeating all of our analysis above for the $c=0$ case with $\gamma^*$ and $\hat \gamma$ replaced with $1-\gamma^*$ and $1-\hat \gamma$, respectively. 
\end{proof}

Combining \cref{prop:c0_lowerbound} and \cref{prop:c1_lowerbound}, we have by our decomposition of $\ldsm$ from the start of this section that 
\[\ldsm (\pdist{\hat \gamma},\pdist{\gamma^*})\geq C (\hat\gamma - \gamma^*)^2\]
where
\[C=\frac{ \epsilon ( 1-\Phi(\tau))}{16\sqrt{\beta_1}\; \hat\gamma (1-\hat\gamma) \gamma^* (1-\gamma^*)}\left[(1-\gamma^*)\frac{\tau\sqrt{\tau^2 + R_0}+\frac{1}{2}\tau^2 + \frac{1}{2}R_0}{\sqrt{\ln(\separation)- \ln\fitparenth{\tau+ \frac{1}{2}\sqrt{\tau^2 + R_0}}}}+\gamma^* \frac{\tau\sqrt{\tau^2 + R_1}+\frac{1}{2}\tau^2 + \frac{1}{2}R_1}{\sqrt{\ln(\separation)- \ln\fitparenth{\tau+ \frac{1}{2}\sqrt{\tau^2 + R_1}}}}\right]\]
for all $\tau\geq \separation\sqrt{\alpha_1}$ and $\epsilon\in [ \epsilon_{\mathrm{min}}, \epsilon_{\mathrm{\max}}]$, where
\begin{align*}
\epsilon_{\mathrm{max}}&=\min\set{\epsilon_{\mathrm{max}}^0,\epsilon_{\mathrm{max}}^1}=\exp\fitparenth{-\frac{1}{8}\fitbracket{(\hat r-r^*)^2 + \max\set{\fitparenth{\hat r+r^*-\tau^2}_+^2,\fitparenth{-(\hat r+r^*)-\tau^2}_+^2}}}
\\
\epsilon_{\mathrm{min}}&=\max\set{\epsilon_{\mathrm{min}}^0,\epsilon_{\mathrm{min}}^1}=\exp\fitparenth{-\frac{1}{8}\fitbracket{\min\set{(\separation(\separation - 2\tau)+(\hat r+r^*))^2,(\separation(\separation - 2\tau)-(\hat r+r^*))^2} + (\hat r-r^*)^2}}
\end{align*}

To simplify this expression, note that if 
\[g(x):=\frac{x(\tau+\frac{1}{2}x)}{\sqrt{\ln(\separation) - \ln(\tau + \frac{1}{2}x)}}\]
then direct calculation shows that if $h(x):=\ln(\separation) - \ln(\tau + \frac{1}{2}x)$,
\[g'(x)=\frac{\frac{x}{4}+(\tau+x)h(x)}{h(x)^{3/2}}>0\]

Therefore, $g$ is increasing, meaning
\begin{align*}
    C&=\frac{ \epsilon ( 1-\Phi(\tau))}{16\sqrt{\beta_1}\; \hat\gamma (1-\hat\gamma) \gamma^* (1-\gamma^*)}\left[(1-\gamma^*)g\fitparenth{\sqrt{\tau^2+R_0}}+\gamma^* g\fitparenth{\sqrt{\tau^2+R_1}}\right]
    \\
    &\geq \frac{ \epsilon ( 1-\Phi(\tau))}{16\sqrt{\beta_1}\; \hat\gamma (1-\hat\gamma) \gamma^* (1-\gamma^*)}g\fitparenth{\min\set{\sqrt{\tau^2+R_0}, \sqrt{\tau^2+R_1}}}
    \\
    &= \frac{ \epsilon ( 1-\Phi(\tau))}{16\sqrt{\beta_1}\; \hat\gamma (1-\hat\gamma) \gamma^* (1-\gamma^*)}g\fitparenth{\sqrt{\tau^2+\sqrt{8\ln \frac{1}{\epsilon}-(\hat r - r^*)^2}- \abs{\hat r-r^*}}}
\end{align*}

Let $R_{\mathrm{min}}:=\sqrt{8\ln \frac{1}{\epsilon}-(\hat r - r^*)^2}- \abs{\hat r-r^*}$. To obtain the bound from the body of the paper, we now specify values for $\tau$ and $\epsilon$ which simplify the expression. First, we will set $\tau$ to be its minimum value: $\tau=\separation\sqrt{\alpha_1}$. We choose this because numerically, it seems to be the value of $\tau$ that maximizes the expression. This gives us
\[C\geq \frac{ \epsilon ( 1-\Phi(\separation\sqrt{\alpha_1}))}{16\sqrt{\beta_1}\; \hat\gamma (1-\hat\gamma) \gamma^* (1-\gamma^*)} \frac{\sqrt{\separation^2\alpha_1 + R_{\mathrm{min}} }\fitparenth{\separation\sqrt{\alpha_1} +\frac{1}{2}\sqrt{\separation^2\alpha_1+ R_{\mathrm{min}}}}}{\sqrt{\log(\separation)-\log\fitparenth{\sqrt{\alpha_1}}-\log\fitparenth{\separation+\frac{1}{2}\sqrt{\separation^2+\frac{R_0}{\alpha_1}}}}}\]

Noting that $\log(\separation)<\log\fitparenth{\separation+\frac{1}{2}\sqrt{\separation^2+\frac{R_0}{\alpha_1}}}$, the denominator of the second term is less than $\sqrt{-\log(\sqrt{\alpha_1})}=\frac{1}{4}\beta_1$ (recall $\alpha_t=\exp\fitparenth{-\int_0^t \beta_s\mathrm{d}s}=\exp\fitparenth{-\frac{1}{2}\beta_1 t^2}$). Additionally, note that the numerator of the second term is an increasing function of $\separation$, meaning we can lower-bound it by setting $\separation=0$. Combining these gives us 
\[C\geq \frac{1-\Phi(\separation\sqrt{\alpha_1})}{8\beta_1\; \hat\gamma (1-\hat\gamma) \gamma^* (1-\gamma^*)} \epsilon \,R_{\mathrm{min}} \]

Finally, direct calculation shows that $\epsilon\, R_{\mathrm{min}}$ is maximized when $\epsilon=\exp\fitparenth{-\frac{1}{8}\fitparenth{(\hat r-r^*)^2 + \delta^2}}$ where $\delta = \frac{|\hat r-r^*|+ \sqrt{(\hat r-r^*)^2+16}}{2}$, giving 
\[\epsilon \, R_{\mathrm{min}} =\frac{\sqrt{(\hat r-r^*)^2+16}-|\hat r-r^*|}{2} \exp\fitparenth{-\frac{1}{8}\fitparenth{(\hat r-r^*)^2 + \delta^2}}\]

Note $\delta\leq \abs{\hat r-r^*}+2$ (because $\sqrt{a^2+b^2}\leq |a|+|b|$), and $\delta^2\leq \abs{\hat r-r^*}^2+2\abs{\hat r-r^*}+4$. Applying this, and the fact that $\hat \gamma(1-\hat \gamma)\leq\frac{1}{4}$
\begin{align*}
C&\geq \frac{\fitparenth{1-\Phi(\separation \sqrt{\alpha_1})}\fitparenth{\sqrt{(\hat r-r^*)^2+16}-|\hat r-r^*|}}{4\beta_1\;  \gamma^* (1-\gamma^*)}\exp\fitparenth{-\frac{1}{4}\fitparenth{(\hat r-r^*)^2 +\abs{\hat r-r^*}} - \frac{1}{2}}
\\
&=\frac{\fitparenth{1-\Phi(\separation \sqrt{\alpha_1})}\fitparenth{\sqrt{(\hat r-r^*)^2+16}-|\hat r-r^*|}}{4\sqrt{e}\beta_1\;  \gamma^* (1-\gamma^*)}\exp\fitparenth{-\frac{1}{4}\fitparenth{(\hat r-r^*)^2 +\abs{\hat r-r^*}}}
\end{align*}

Finally, for $\hat \gamma\in [a,1-a]$, $\abs{\hat r - r^*} \leq \log\frac{1-a}{a}+\abs{\log \frac{1-\gamma^*}{\gamma^*}}$. Our expression is a decreasing function of $\abs{\hat r-r^*}$, and thus for all $\hat \gamma\in [a,1-a]$, if we let $k:=\log\frac{1-a}{a}+\abs{\log \frac{1-\gamma^*}{\gamma^*}}$
\[C\geq \frac{\fitparenth{1-\Phi(\separation \sqrt{\alpha_1})}\fitparenth{\sqrt{k^2+16}-k}}{7\beta_1\;  \gamma^* (1-\gamma^*)}\exp\fitparenth{-\fitparenth{k^2 +k}/4}\]
where we additionally used $4\sqrt{e}<7$. Bringing this all together, 
\begin{align}
   \min_{\hat \gamma \in [a,1-a]}\frac{\ldsm(\pdist{\hat\gamma},\pdist{\gamma^*})}{(\hat\gamma-\gamma^*)^2}>\frac{(1-\Phi(\norm{\mu_1-\mu_0}\sqrt{\alpha_1}))\fitparenth{\sqrt{k^2+16}-k}}{7\beta_1\gamma^*(1-\gamma^*)}\exp\fitparenth{-(k^2+k)/4} 
\end{align}
\qed


\section{Additional numerical experiments}\label{sec:Additional Numerics}

\subsection{Noise schedules and time reparameterization}\label{sec: Reparam}
In this work, we reference two noise schedules used in practice. The linear noise schedule is defined by $\beta_t:=\beta_1 t$ for a constant $\beta_1>0$, such that $\alpha_t=\exp\fitparenth{-\frac{\beta_1}{2}t^2}$. The squared-cosine noise schedule \citep{hang2024improvednoiseschedulediffusion}

$    \bar\alpha_t=\cos\fitparenth{\frac{t+s}{1+s}\cdot \frac{\pi}{2}}^2$ and $\alpha_t =\frac{\bar\alpha_t}{\bar\alpha_0}$
for $s=0.008$. 


We want to analyze the effect of artificially reducing the sensitivity index on the estimation of $\gamma^*$. To do this, we can reparameterize the time in the forward process as $t=t_{a,b,v}(\tau)$, where 
\[t_{a,b,v}(\tau)=\begin{cases}
    \tau & 0\leq \tau<a
    \\
    a + v(\tau-a)& a \leq \tau <a + \frac{b-a}{v}
    \\
    b + \frac{1-b}{1-\fitparenth{a + \frac{b-a}{v}}}\fitparenth{t - a - \frac{b-a}{v}} & a + \frac{b-a}{v} \leq t\leq 1
\end{cases}\]

Under this parameterization, if we let $\vect y_\tau:= \vx_{t(\tau)}$ where $\vx_t$ follows the Variance Preserving SDE with noise schedule $\set{\alpha_t}$, then $\vect y_t$ satisfies
\begin{align}
    \mathrm{d}\vect y_\tau &= -\frac{\beta_{t(\tau)}}{2}\vect y_\tau \frac{\mathrm{d}t}{\mathrm{d} \tau}\mathrm{d}\tau + \sqrt{\beta_{t(\tau)}}\cdot \sqrt{\frac{\mathrm{d}t}{\mathrm{d}\tau}}\mathrm dW_t
    \\
    &= -\frac{\tilde{\beta}_{\tau}}{2}\vect y_\tau \mathrm{d}\tau + \sqrt{\tilde{\beta}_{\tau}}\mathrm{d}W_t
\end{align}
where $\tilde{\beta}_\tau:= \beta_{t(\tau)} \cdot \frac{\mathrm{d}t}{\mathrm{d}\tau}$. In particular, reparameterizing time is equivalent to choosing a different noise schedule $\set{\tilde{\alpha}_\tau}$. For $v\geq 1$, $t_{a,b,v}$ maps the interval $[a,b]$ to $\fitbracket{a, a+\frac{b-a}{v}}$, scaling the length of the interval by $1/v$. To reduce the value of $L(\gamma^*)$, we let $[a,b]$ contain the region of time where $\frac{\lsm(\pdist[t]{\gamma}, \pdist[t]{\gamma^*})}{(\gamma-\gamma^*)^2}>\epsilon$ for fixed $\epsilon>0$ and increase the value of $v$.

\subsection{Neural network experiment details}

For the score models $\vs_\btheta$ in our experiments, we use deep ResNets \citep{He2015DeepRL} with width equal to the ambient dimension and depth $8$. Unless otherwise specified, we use a linear noise schedule with $\beta_1=20$ and a time discretization of $T=4000$ steps on the interval $[10^{-5},1]$ for training and $[10^{-3},1]$ for sampling, following \citep{song2021scorebasedgenerativemodelingstochastic}.

We train our models using the loss described in this paper (and \citep{song2021scorebasedgenerativemodelingstochastic}) using AdaBelief \citep{zhuang2020adabelief} with a linear warmup and cosine decay learning rate. We use JAX \citep{jax2018github} and Equinox \citep{kidger2021equinox} for our neural network code, SciPy \citep{2020SciPy-NMeth} for additional machine learning algorithms, Matplotlib \citep{Hunter2007matplotlib} for figure generation, and Weights and Biases \citep{wandb} for experiment tracking.
\footnote{We will make our code available upon publication.}

\subsection{MNIST ablation study details}\label{sec: MNIST Exp Details}

The architecture of our autoencoder is simple: two convolutional layers followed by a dense layer for the encoder, and the reverse for the decoder. For our MNIST experiments, we use a latent dimension of $12$. We train the autoencoder on all examples of $1$s and $8$s in MNIST. We additionally train an MNIST classifier on only $1$s and $8$s using the architecture in the Equinox documentation \citep{kidger2021equinox}.

Unlike in the case of Gaussian mixture models, we do not know the analytical score function for the latents of our autoencoder, and thus estimation of $\ldsm(\pdist[t]{\gamma}, \pdist[t]{\gamma^*})$ and $L(\gamma^*)$ analytically is intractable. Instead, for $\gamma^*\in \set{0.2,0.22,...,0.8}$, we train diffusion models $\vs_{\btheta_\gamma}$ on data from $\pdist{\gamma}$ as approximations for the analytical scores using a linear noise schedule. We keep the number of training examples constant across all experiments (5800, the smaller of the counts of 1s and 8s in MNIST) and train the models identically. For each $t\in [0,1]$, we then have the estimate \[\ldsm(\pdist[t]{\gamma}, \pdist[t]{\gamma^*})\approx \expect[\vx_t\sim \pdist[t]{\gamma^{*}}]{\lambda_t \norm{\vs_{\btheta_{\gamma^*}}(\vx_t, t) - \vs_{\btheta_{\gamma'}}(\vx_t, t)}^2}\]

We generate all samples used in \Cref{fig: MNIST experiment} using $\vs_{\btheta_{0.6}}$, varying the noise schedule used for sampling.

% \subsection{High-dimensional GMM experiments}\label{sec:High Dim GMM Exp}

% To demonstrate that our results hold in higher dimensions, we train diffusion models on $\pdist{\gamma}=(1-\gamma) \;\mathcal{N}(\vect 0, \vect I)+\gamma \;\mathcal{N}\fitparenth{\frac{12}{\sqrt{d}}\vect 1, \vect I}$ in dimension $d=100$ for $\gamma^*\in\set{0.05,0.06,...,0.5}$ fixing the number of training samples at $10^{15}$. We then estimate the mixture weights, means, and covariances using the EM algorithm. Due to the computational cost of running these experiments, we run only one trial for each $\gamma^*$ and therefore do not report error bars.


% \begin{figure}[ht]
%     \centering
%     \begin{subfigure}[htbp]{0.19\linewidth}
%         \includegraphics[width=\columnwidth]{images/figure_one_a.pdf}
%         \caption{}
%         \label{fig: high dim gmm mix}
%     \end{subfigure}
%     \begin{subfigure}[htbp]{0.19\linewidth}
%         \includegraphics[width=\columnwidth]{images/figure_one_d.pdf}
%         \caption{}
%         \label{fig: high dim gmm means}
%     \end{subfigure}
%     \begin{subfigure}[htbp]{0.19\linewidth}
%         \includegraphics[width=\columnwidth]{images/figure_one_d.pdf}
%         \caption{}
%         \label{fig: high dim gmm covs}
%     \end{subfigure}
%     \caption{Estimation errors in }
%     \label{fig:Langevin_vs_Diffusion}
% \end{figure}

\subsection{DSSI for other parameters }\label{sec: Other Parameters}

Here, we apply our DSSI framework to parameters other than the mixture weight in a Gaussian model. In what follows, we will let $\Sigma_0$ and $\Sigma_1$ be fixed random positive definite matrices, such that 
\[\Sigma_i = A^\top A+10^{-8}I, \;\;A_{ij}\sim \mathcal{N}(0,1)\]
for $i\in \set{0,1}$ and use a linear noise schedule. We focus here on one-dimensional families of distributions for consistency with the body of the paper, although we note that our framework can be extended to multi-parameter recovery directly. For this section, we examine GMMs in dimension $d=10$.

For our analysis of the sensitivity of $\ldsm$ to the modes of the mixture components, we consider the one-dimensional family of distributions 
\[\pdist{\delta}:=0.6 \; \mathcal{N}(\vect 0, \vect \Sigma_0) + 0.4 \; \mathcal{N}\fitparenth{\frac{12\;\delta}{\sqrt{d}}\vect 1,\vect\Sigma_1}\] 
in dimension $d=10$ where we range $\delta$ between $0$ and $1$. \Cref{fig: means heatmap} shows $\lsm(\pdist[t]{\delta}, \pdist[t]{\delta^*})$ as a function of $\delta$ and $t$, and \Cref{fig: means index} shows $\frac{\ldsm(\pdist[t]{\delta}, \pdist[t]{\delta^*})}{(\delta - \delta^*)^2}$ as a function of $\delta$ for multiple values of $\delta^*$. We see that, as with the mixture weight, $\lsm(\pdist[t]{\delta}, \pdist[t]{\delta^*})$ is nonnegligible only for intermediate times in the forward process. However, note that the minimum value of $L(\delta^*)$ we see here is greater than $2$ (compared to $0.37$ in \Cref{fig: gmm score sensitivity}), meaning that our theory implies a $5.4\times$ stronger control on the error in estimation $\mu$ using a diffusion model trained to the same level of error under $\ldsm$.

\begin{figure}
    \centering
    \begin{subfigure}[htbp]{0.45\linewidth}
        \includegraphics[width=\columnwidth]{images/mean_hmap.pdf}
        \caption{}
        \label{fig: means heatmap}
    \end{subfigure}
    \hfill
    \begin{subfigure}[htbp]{0.45\linewidth}
        \includegraphics[width=\columnwidth]{images/means_rat.pdf}
        \caption{}
        \label{fig: means index}
    \end{subfigure}
    \caption{Sensitivity of $\ldsm$ to means of mixture components. (a) $\lsm(\pdist[t]{\delta}, \pdist[t]{\delta^*})$ as a function of $t$ and $\delta$. (b)  $\frac{\ldsm(\pdist{\delta}, \pdist{\delta^*})}{(\delta - \delta^*)^2}$ as a function of $\delta$.}
\end{figure}



We now apply our DSSI framework to covariance matrix recovery. Consider the one-dimensional family of distributions 
\[\pdist{\kappa}:=0.6 \; \mathcal{N}(\vect 0, \vect I) + 0.4 \; \mathcal{N}\fitparenth{\frac{9.0}{\sqrt{d}}\vect 1,\vect\Sigma_\kappa}\] 
in dimension $d=10$ where we range $\delta$ between $0$ and $1$. Here, we define 
\[\vect \Sigma_\kappa:= \vect \Sigma_0^{1/2}\fitparenth{\vect \Sigma_0^{-1/2}\vect \Sigma_1\vect \Sigma_0^{-1/2}}^{\kappa}\vect \Sigma_0^{1/2}\]
to be the geodesic interpolating between $\vect \Sigma_0$ and $\vect \Sigma_1$.  \Cref{fig: variance heatmap} shows $\lsm(\pdist[t]{\kappa}, \pdist[t]{\kappa^*})$ as a function of $\kappa$ and $t$, and \Cref{fig: variance index} shows $\frac{\ldsm(\pdist{\kappa}, \pdist{\kappa^*})}{(\kappa - \kappa^*)^2}$ as a function of $\kappa$.

Again, $\lsm(\pdist[t]{\delta}, \pdist[t]{\delta^*})$ is nonnegligable for intermediate times in the forward process. We additionally see that $L(\kappa^*)$ is greater than $0.6$ uniformly across $\kappa$. 


\begin{figure}
    \centering
    \begin{subfigure}[htbp]{0.45\linewidth}
        \includegraphics[width=\columnwidth]{images/variance_hmap.pdf}
        \caption{}
        \label{fig: variance heatmap}
    \end{subfigure}
    \hfill
    \begin{subfigure}[htbp]{0.45\linewidth}
        \includegraphics[width=\columnwidth]{images/variance_rat.pdf}
        \caption{}
        \label{fig: variance index}
    \end{subfigure}
    \caption{Sensitivity of $\ldsm$ to covariance matrix of mixture components. (a) $\lsm(\pdist[t]{\kappa}, \pdist[t]{\kappa^*})$ as a function of $t$ and $\kappa$. (b)  $\frac{\ldsm(\pdist{\kappa}, \pdist{\kappa^*})}{(\kappa - \kappa^*)^2}$ as a function of $\kappa$.}
\end{figure}



% \newpage
% \input{checklist.tex}
\end{appendices}
\end{document}
